% Elsevier Editorial Manager (EM) build target -- pdflatex-compatible.
%
% EM runs TeX Live 2022 with pdflatex by default and often ignores the
% `%!TEX TS-program = xelatex` directive, so this wrapper deliberately
% avoids XeTeX-only packages (fontspec, unicode-math). Table column
% widths use `calc` + `\linewidth` which are engine-agnostic. The
% companion `main.tex` (in the same directory) uses XeLaTeX for the
% local publication-quality build; both compile the same body.

\documentclass[preprint,11pt,authoryear,a4paper]{elsarticle}

\usepackage[utf8]{inputenc}
\usepackage[T1]{fontenc}
\usepackage{lmodern}       % Latin Modern -- same metrics as Computer Modern

\usepackage{amsmath,amssymb,amsthm}
\usepackage{booktabs}
\usepackage{graphicx}
\usepackage{siunitx}
\usepackage{xurl}
\usepackage{hyperref}
\hypersetup{hidelinks}
\usepackage[capitalise,nameinlink]{cleveref}
\usepackage{xcolor}
\usepackage{longtable}
\usepackage{enumitem}
\usepackage{array}
\usepackage{calc}
\usepackage{multirow}

\usepackage[a4paper,margin=2.5cm]{geometry}
\usepackage{microtype}
\usepackage[font=small,labelfont=bf,labelsep=period]{caption}
\setlength{\parindent}{1.5em}

% Map the handful of glyphs that survive postprocess_tex.py's
% ensuremath step. Use \DeclareUnicodeCharacter (inputenc's native API)
% rather than \newunicodechar so we don't emit "Redefining Unicode
% character" warnings for §, –, †, ‡ that inputenc already handles.
\DeclareUnicodeCharacter{25AB}{\ensuremath{\,\square\,}}  % ▫
\DeclareUnicodeCharacter{03C1}{\ensuremath{\rho}}         % ρ
\DeclareUnicodeCharacter{0394}{\ensuremath{\Delta}}       % Δ
\DeclareUnicodeCharacter{00A7}{\S}                        % §
\DeclareUnicodeCharacter{2013}{--}                        % –
\DeclareUnicodeCharacter{2020}{\textsuperscript{\dag}}    % †
\DeclareUnicodeCharacter{2021}{\textsuperscript{\ddag}}   % ‡

% Pandoc compatibility shims.
\providecommand{\tightlist}{\setlength{\itemsep}{0pt}\setlength{\parskip}{0pt}}
\providecommand{\pandocbounded}[1]{#1}
\providecommand{\textsubscript}[1]{\ensuremath{_{\text{#1}}}}
\providecommand{\real}[1]{#1}

% Corrected 2026-08-15: production reported the source PDF had no
% figures. The prior \emfig macro rendered a placeholder box naming a
% TIFF file rather than embedding the figure, on the assumption that
% Elsevier's typesetter would insert the TIFF at production. That
% assumption produced a review PDF with no visible figures. This
% version embeds the PNG directly via \includegraphics so pdflatex
% produces a fully self-contained PDF with real figures.
\newcommand{\emfig}[1]{%
  \includegraphics[width=0.92\linewidth]{#1.png}%
}
\newcommand{\emfigOLD}[1]{%
  \fbox{\parbox[c][4.5cm][c]{0.75\linewidth}{%
    \centering
    \textit{Figure to be inserted at production:}\\[6pt]%
    \texttt{#1.tif}%
  }}%
}

\journal{BioSystems}

\begin{document}

\begin{frontmatter}

\title{Robust error-minimization in the genetic code across physicochemical metrics and variant codes: a graph-theoretic analysis in $\mathrm{GF}(2)^6$}

\author[lc]{Paul Clayworth\fnref{eq}}
\author[bs]{Sergey Kornilov\corref{cor}\fnref{eq}}

\affiliation[lc]{organization={Logocentricity Inc.}, addressline={Austin, TX}, country={USA}}
\affiliation[bs]{organization={Biostochastics, LLC}, addressline={Seattle, WA}, country={USA}}

\cortext[cor]{Corresponding author. \texttt{sergey@biostochastics.com}}
\fntext[eq]{Equal contribution.}

\begin{abstract}
The standard genetic code reduces the impact of point mutations, but the robustness of this property across physicochemical metrics, naturally variant codes, and codon-reassignment mechanisms remains incompletely quantified. Embedding the 64 codons in $\mathrm{GF}(2)^6$ represents the hypercube $Q_6$ as a coordinate-dependent subgraph of the encoding-independent single-nucleotide mutation graph $H(3,4)$, and enables continuous $\rho$-interpolation between the two. Under a block-preserving null ($n = 10{,}000$), the standard code is significantly low-cost across four established, code-independent physicochemical distance metrics with partially overlapping content (Grantham $p = 0.006$; Miyata $p < 0.001$; Woese polar requirement $p = 0.003$; Kyte--Doolittle hydropathy $p = 0.001$), and the signal strengthens monotonically as $\rho$ moves $Q_6 \to H(3,4)$. A structure-aware sensitivity analysis under the alignment-derived ProtSub matrix (Jia \& Jernigan 2021) yields the most extreme percentile of any measure tested ($p = 0.0004$; all five p-values pass Bonferroni at $\alpha = 0.05$). Across the 27 NCBI translation tables, near-optimality is preserved: 11 of 12 informative-distance variants retain top-5\% placement after BH--FDR correction. Natural codon reassignments avoid disrupting codon-family connectivity: under $H(3,4)$, observed events are topology-breaking at relative risk $0.32$ versus the candidate landscape (permutation $p \leq 10^{-4}$), robust to alternative topology definitions and base-to-bit encodings. Event-level conditional-logit modelling shows that topology avoidance and local physicochemical cost provide complementary, only weakly correlated signal ($r_s = 0.15$), and that topology adds explanatory value beyond physicochemistry under both $Q_6$ and encoding-independent $H(3,4)$ adjacency. Retrospective reanalysis of nine genome-recoding datasets is consistent with codon-family topology operating as an evolutionary-trajectory constraint distinct from acute engineering fitness. The contribution is the second axis: code evolution is jointly constrained by physicochemical smoothness and codon-family topological integrity, and these two constraints are partly independent.
\end{abstract}

\begin{keyword}
genetic code evolution \sep error minimization \sep $\mathrm{GF}(2)^6$ representation \sep codon reassignment \sep physicochemical distance \sep codon-family topology \sep conditional logit
\end{keyword}

\begin{highlights}
\item Codon-space geometry links genetic-code robustness and reassignment paths
\item Standard and variant codes preserve broad physicochemical error minimization
\item Reassignments are depleted for codon-family topology-breaking moves
\item Conditional-logit models separate topology from physicochemical similarity
\item Synthetic recoding shows boundary conditions for natural-code constraints
\end{highlights}

\end{frontmatter}

\section{Introduction}\label{sec:intro}

The standard genetic code maps 61 sense codons to 20 amino acids through
a pattern that has long been recognized as non-random. \citep{woese1965}
first noted that similar codons tend to encode amino acids with similar
physicochemical properties, and \citep{freeland1998} demonstrated
quantitatively that the standard code sits among approximately 1 in
\(10^{6}\) random codes for mutational error minimization under the
Woese polar requirement distance. These findings established that the
code's structure reduces the fitness impact of point mutations, but
whether this optimality is specific to a single metric or robust across
distinct physicochemical parameterizations has not been systematically
tested.

Every codon comprises three nucleotides drawn from
\(\left\{ C,U,A,G \right\}\). By choosing a bijection
\(\varphi:\left\{ C,U,A,G \right\} \rightarrow {\text{ GF}(2)}^{2}\)
(for instance, \(C \mapsto (0,0)\), \(U \mapsto (0,1)\),
\(A \mapsto (1,0)\), \(G \mapsto (1,1)\)), each codon maps to a vertex
of the 6-dimensional binary hypercube \(Q_{6} = {\text{ GF}(2)}^{6}\).
The genetic code then becomes a \emph{coloring} of \(Q_{6}\) by 21
labels (20 amino acids plus the stop signal), and single-nucleotide
mutations correspond to edges of \(Q_{6}\) or, more precisely, to a
subgraph of the complete mutation graph \(H(3,4)\).

This representation is not new in principle: binary encodings of the
genetic code appear in the mathematical biology literature (e.g.,
\citep{antoneli2011}). Its value lies not in the encoding itself but in
the analytical decomposition it enables: the complete single-nucleotide
mutation graph \(H(3,4)\) (288 edges) separates into 192
Hamming-distance-1 edges (single-bit changes) and 96 within-nucleotide
distance-2 edges (both bits of one nucleotide position flip
simultaneously), allowing systematic interpolation via a weight
parameter \(\rho \in \lbrack 0,1\rbrack\). This decomposition does not
correspond to the biological transition/transversion partition: under 16
of the 24 possible 2-bit encodings, the Hamming-1 edges contain an equal
mixture of transitions and transversions per nucleotide position (2 of
each among 4 Hamming-1 pairs), while the remaining 8 encodings place
both transitions on diagonal (Hamming-2) edges. The \(\rho\) parameter
should therefore be interpreted as a diagonal-edge inclusion weight, not
a transition/transversion weight. Previous work has not exploited this
decomposition to test error-minimization across multiple physicochemical
metrics, nor extended the analysis to variant genetic codes. Doing so
lets us examine two questions that single-axis analyses cannot. First,
whether error-minimization is a property of the code itself rather than
of any particular distance function. Because the established
physicochemical measures (composition--polarity--volume,
polarity--volume, chromatographic polar requirement, hydropathy) are
partially overlapping, we do not treat concordance across them as
independent replication, but as a convergent sensitivity envelope: if
the standard code lies in the low-cost tail under several established
but non-identical parameterizations, the result is not confined to a
single chosen scale (overlap quantified via pairwise and partial
Spearman correlations in §3.1). Second, whether codon-family
\emph{connectivity} acts as a second, partly independent axis of
constraint on natural reassignment events, distinct from physicochemical
cost. If so, code evolution is bounded not only by which codes are
low-cost but by which trajectories through code-space leave the decoding
substrate intact, a constraint on transitions rather than on states.
These two motivating questions are operationalized as three concrete
tests:

\begin{enumerate}
\item
  \textbf{Is the standard code optimal under multiple physicochemical
  metrics?} We test whether the standard code's edge-mismatch score is
  extreme relative to quartet-pattern shuffle null models across four
  established, code-independent physicochemical distance measures with
  partially overlapping content (Grantham, Miyata, Woese polar
  requirement, and Kyte--Doolittle hydropathy), extending
  \citep{freeland1998} from a single metric to a cross-metric
  sensitivity envelope.
\item
  \textbf{Is this structure preserved by evolution?} We ask whether
  variant genetic codes (NCBI translation tables 2--33) maintain
  error-minimization, and whether natural codon reassignment events
  preferentially avoid disrupting the topological connectivity of amino
  acid codon families.
\item
  \textbf{What are the genomic correlates of disruption?} We test
  whether organisms whose variant codes break the connectivity of an
  amino acid's codon graph show elevated tRNA gene copy numbers for the
  affected amino acid, using tRNAscan-SE--verified data from 18 genomes
  spanning 5 variant genetic codes across Alveolata, Opisthokonta,
  Excavata, and Mollicutes; 15 of these 18 genomes populate the
  24-pairing enrichment analysis, with the remaining 3 (Blastocrithidia
  and two Mycoplasmoides species) discussed as mechanistic boundary
  cases in §4.3 because their reassignment routes act without changing
  tRNA gene counts.
\end{enumerate}

Each of these three tests carries an explicit decision criterion. For
(1), a null result would be quartet-pattern shuffle null \(p > 0.05\) on
any of the four physicochemical metrics; discordance among metrics (some
significant, some not) would restrict any claim of optimality to a
specific parameterization rather than to the code itself. For (2), a
null result would be either failure of BH-corrected per-table
significance in the majority of informative-distance tables
(\(d_{H} \geq 3\) reassignments from the standard code), or an observed
topology-breaking rate not distinguishable from the candidate-landscape
rate (hypergeometric \(p > 0.05\) or risk-ratio 95\% CI including 1.0).
For (3), the originally specified decision gate was the
\textbf{worst-case} Stouffer \(p\)-value across all maximal independent
pairing subsets; that gate was not met (\(p = 0.104\);
\Cref{sec:res-trna}), so the tRNA result is reported as
\textbf{exploratory} rather than as confirmatory support for the
compensation-via-tRNA-duplication hypothesis. We report the median and
full MIS distribution descriptively (median \(p = 0.037\); 264 of 332
MIS below 0.05); these motivate an exploratory decoding-accommodation
signal rather than a confirmatory enrichment claim. The pre-specified
topology-breaking-only subset (\(n = 4\)) is examined as a complementary
mechanistic test and is also null (Stouffer \(p = 0.387\)). We report
all three tests with the corresponding numbers regardless of outcome.

To delimit the framework's scope, we also test four conjectural
extensions previously associated with this project (one of us, P.C.;
\citep{clayworth2026}): a Serine distance-4 invariance claim, PSL(2,7)
symmetry, a holomorphic embedding extending a character of
\({\text{GF}(8)}^{\ast}\), and a KRAS--Fano clinical prediction. These
tests distinguish supported graph-theoretic results from encoding
artifacts and unsupported algebraic extensions
(\Cref{sec:res-negatives}; KRAS--Fano detail in Supplement §S13).

The paper is organized as follows. Section 2 describes the encoding
formalism, graph decomposition, null models, and statistical methods,
including a retrospective cross-study reanalysis of nine published
genome-recoding datasets (eight analyzed quantitatively). Section 3
presents the four supported findings (cross-metric coloring optimality,
per-table preservation, \(\rho\)-robustness, and topology avoidance),
the exploratory tRNA enrichment result, the synthetic-recoding
reanalysis, and additional exploratory observations. Section 4 discusses
the graph-theoretic interpretation, the relationship to frozen-accident
versus adaptive hypotheses \citep{koonin2009}, and an exploratory
three-layer interpretation. All analyses are reproducible via the
open-source \texttt{codon-topo} pipeline (version 0.6.1, seed 135325).

\section{Methods}\label{sec:methods}

\subsection{Binary encoding of the genetic code}\label{sec:encoding}

We encode each nucleotide base as a 2-bit vector via the default
bijection \(\varphi:C \mapsto (0,0)\), \(U \mapsto (0,1)\),
\(A \mapsto (1,0)\), \(G \mapsto (1,1)\). A codon \(b_{1}b_{2}b_{3}\) is
then represented as the concatenation
\(\varphi(b_{1})\|\varphi(b_{2})\|\varphi(b_{3}) \in {\text{ GF}(2)}^{6}\),
and the 64 codons become the 64 vertices of the 6-dimensional hypercube
\(Q_{6}\). This default bijection was adopted in companion
methodological work (\citep{clayworth2026}) because it places the
standard code's nine 2-fold-degenerate amino acids on bit-5 differences.
The visualization-clarity rationale, however, is mildly circular: the
bit-5 two-fold filtration is itself classified as Tautological in the
present claim hierarchy (Supplement §S1), so the encoding was chosen to
make a tautological property visible. We therefore use the default
encoding only as an analytical convenience and make no claim that it is
biologically privileged. Two specific encoding-dependent results are
reported under all 24 base-to-bit bijections: (i) the
coloring-optimality Grantham quartet-pattern shuffle null (Supplement
§S2), and (ii) the \(Q_{6}\) topology-avoidance landscape depletion
(Supplement §S4). The primary topology-avoidance test uses the
encoding-independent \(H(3,4)\) adjacency (§2.3.4). The
conditional-logit models M1--M4 are not swept across encodings; their
\(H(3,4)\) verification variants replace the \(\Delta_{topo,Q\_ 6}\)
topology feature with the encoding-independent \(\Delta_{topo,H(3,4)}\)
feature but retain a local physicochemical feature evaluated on
\(Q_{6}\) Hamming-1 edges, so those variants are best described as
\textbf{topology-verification} rather than as fully encoding-independent
combined models.

Under this encoding, two codons that differ by a single-nucleotide
substitution in which only one bit of the 2-bit pair changes are
adjacent in \(Q_{6}\) (Hamming distance 1). However, transversions that
flip both bits within a nucleotide position correspond to Hamming
distance 2. The full single-nucleotide mutation graph is therefore the
Hamming graph \(H(3,4) = K_{4}\, \ensuremath{\,\square\,}\, K_{4}\, \ensuremath{\,\square\,}\, K_{4}\) (the Cartesian
product of three complete graphs on the four nucleotide states;
equivalently, \(K_{4}^{\ensuremath{\,\square\,}3}\)), with 64 vertices, regular degree 9, and
288 undirected edges. It contains \(Q_{6}\) as a 192-edge subgraph:
\(Q_{6}\) contributes the 192 Hamming-distance-1 edges (single-bit
changes within one 2-bit nucleotide), while the remaining 96
within-nucleotide diagonal (Hamming-distance-2) edges complete
\(H(3,4)\). To address the concern that \(Q_{6}\) misses approximately
one-third of single-nucleotide mutations, we introduce the weighted
score \(F_{\rho}\) that interpolates between pure \(Q_{6}\)
(\(\rho = 0\)) and full \(H(3,4)\) (\(\rho = 1\)); see Section 2.3.
Throughout this paper \(H(3,4)\) denotes this nucleotide-level mutation
graph, which is encoding-independent (every two-bit bijection from
\(\left\{ A,C,G,U \right\}\) to \(\left\{ 0,1 \right\}^{2}\) yields the
same \(H(3,4)\)); \(Q_{6}\) is encoding-dependent, since the partition
into Hamming-1 vs Hamming-2 edges depends on the bijection.

There are \(4! = 24\) such bijections. All encoding-dependent results
are tested across all 24; encoding-invariant properties (such as
Serine's disconnection at \(\varepsilon = 1\), where \(\varepsilon\)
denotes the Hamming distance threshold in \({\text{GF}(2)}^{6}\)) are
noted as such. Coloring optimality is significant (\(p < 0.05\)) under
every encoding (full sweep in the Supplementary Material).

\subsection{Edge-mismatch objective function}\label{sec:objective}

The genetic code assigns each vertex \(v \in Q_{6}\) a label
\(c(v) \in \mathcal{A}\), where \(\mathcal{A}\) comprises the 20 amino
acids and the stop signal. The edge-mismatch score is

\[F(c) = \sum_{\left\{ v,w \right\}:d(v,w) = 1}\Delta(c(v),c(w))\]
\protect\phantomsection\label{eq:mismatch}{}

where the sum ranges over all 192 edges of \(Q_{6}\) (pairs of vertices
at Hamming distance 1), and \(\Delta\) is a physicochemical distance
between amino acids. We test four established, code-independent
physicochemical distance measures with partially overlapping content:
the \citep{grantham1974} composite distance (composition, polarity,
volume; range 5--215), the \citep{miyata1979} normalized Euclidean
distance (polarity and volume only; range 0.06--5.13), the Woese polar
requirement absolute difference (range 0--8.2;
\citep{woese1966, woese1973}; used by \citep{freeland1998}), and the
\citep{kyte1982} hydropathy absolute difference (range 0--9.0; used by
\citep{haig1991}). All four are \textbf{code-independent}
physicochemical scales, derived from properties of amino acids in
isolation (chromatographic partition coefficient, composition, polarity,
volume, hydropathy) rather than from observed substitution patterns in
coded proteins; \citep{kosiol2004} review criteria for principled
amino-acid classification schemes in this same code-independent regime.
This design avoids the tautology that \citep{digiulio2001} identified in
PAM-derived measurements of code optimality, where the substitution
matrix itself reflects the genetic code structure under study;
\citep{koonin2009} explicitly endorse code-independent physicochemical
scales for exactly this reason. As a structure-aware sensitivity
analysis, we also report results under the MSA-derived ProtSub matrix
(\citep{jia2021} 2,320 Pfam domains), which is alignment-derived and
therefore code-dependent in the sense of \citep{digiulio2001}. ProtSub
is treated as a robustness check rather than as primary evidence for
code-origin optimality. Synonymous edges contribute 0; edges involving a
stop codon receive a fixed penalty scaled proportionally to each
metric's maximum. A lower \(F\) indicates a more error-minimizing code.

Stop codons are held fixed across all null models (Section 2.3) because
their assignment is constrained by release-factor recognition geometry
(eRF1 binding UAR/UGA) rather than tRNA decoding, and permuting stops
would conflate two evolutionarily separate optimization problems. The
stop-codon contribution is therefore a constant offset; sensitivity
analysis across penalty values (0, 150, 215, 300) confirms this is
immaterial to the ranking.

\subsection{Null models}\label{sec:nullmodels}

\subsubsection{Quartet-pattern shuffle null}\label{sec:null-fh}

The randomisation ensemble used throughout the coloring-optimality tests
is a \textbf{quartet-pattern shuffle}: we group the 64 codons into 16
blocks of 4, defined by shared first-two-base prefix (the codon
``quartets''). Each quartet's internal 4-tuple of amino-acid labels is
held as an atomic pattern, and the 16 patterns are permuted uniformly at
random among the 16 quartet slots. Quartets containing stop codons are
held fixed. This null preserves both the synonymous codon contiguity
(wobble degeneracy) and each quartet's internal split structure (whether
it is a \((4)\), \((2,2)\), \((3,1)\), or singleton block).

The reassignment biology literature also uses a related null often cited
to \citep{haig1991} and \citep{freeland1998}: it fixes the standard
code's 20 sense-codon families (the unlabeled partition of the 61 sense
codons into synonymous groups, holding stop positions fixed) and
permutes the 20 amino-acid labels uniformly across those families. This
classical Haig--Hurst / Freeland--Hurst amino-acid-permutation null is a
different randomisation ensemble from the quartet-pattern shuffle we
implement. The two ensembles preserve non-nested structural properties:
the quartet-pattern shuffle moves the standard code's atomic 4-codon
quartet patterns among the 14 non-stop quartet slots, so each named
amino acid's total codon count is preserved but which specific codons
decode as that amino acid is not; the AA-permutation null fixes the
codon-family partition but allows a named amino acid's codon count to
change (e.g., Trp's single-codon family may become a 6-codon family if
its label is swapped with Leu's). The two ensembles admit different
numbers of alternative codes (\(\approx 14! \approx 8.7 \times 10^{10}\)
for the quartet-pattern shuffle versus
\(\approx 20! \approx 2.4 \times 10^{18}\) for the AA-permutation null)
and can produce different null tails. We report the four physicochemical
metrics under \textbf{both} nulls in Supplement §S2.1: treating the
eight reported \(p\)-values (four metrics \ensuremath{\times} two nulls, \(n =\)10,000
draws each, seed 135325) as one test family, all eight pass Bonferroni
at \(\alpha = 0.05/8 = 0.00625\). For three of the four code-independent
metrics the classical Haig--Hurst null gives a \textbf{more} extreme
\(p\) than the quartet-pattern shuffle, and for Kyte--Doolittle the
direction reverses. We therefore do not describe either null as
uniformly ``more stringent'' and treat the classical null as a
same-direction sensitivity check on the quartet-pattern-shuffle headline
values.

For the standard code, we drew \(n =\)10,000 null samples (seed 135325)
from the quartet-pattern shuffle and computed the conservative
\(p\)-value as \(\frac{k + 1}{n + 1}\), where \(k\) is the number of
null scores below the observed \(F\).

\subsubsection{Per-table null}\label{sec:null-pertable}

The NCBI Genetic Codes registry
(\url{https://www.ncbi.nlm.nih.gov/Taxonomy/Utils/wprintgc.cgi;} gc.prt
v4.6, retrieved 2026-04-25) currently lists 27 translation tables: codes
1--6, 9--16, and 21--33 (codes 7, 8, and 17--20 are deprecated and
absent from the current registry). All 27 are analyzed in this work. Two
pairs of tables share identical sense-codon mappings and differ only in
their start-codon assignments: tables 1 (Standard) and 11 (Bacterial /
Archaeal / Plant Plastid), and tables 27 (Karyorelict Nuclear) and 28
(Condylostoma Nuclear). Both pairs are retained as separate entries to
match NCBI numbering, but produce identical results for analyses that
depend only on codon\(\rightarrow\)amino-acid mappings; the 27 NCBI
tables thus correspond to 25 distinct sense-codon colorings. Each table
was tested against its own quartet-pattern shuffle null (\(n =\)10,000
per table, common-seed design with seeds = base seed + table ID).
\(P\)-values were corrected for multiple comparisons using the
\citep{benjamini1995} false discovery rate (FDR) procedure. For NCBI
tables with dual-function stop/sense codons (e.g., table 27 lists UGA as
``Stop or Trp''; table 28 lists UAA, UAG as ``Gln or Stop'' and UGA as
``Trp or Stop``), the primary codon\(\rightarrow\)amino-acid analyses
use the amino-acid label given in the NCBI AAs row of the gc.prt
definition; stop functionality is handled only in analyses explicitly
involving stop labels.

\subsubsection{\texorpdfstring{Graph family \(G_{\rho}\) and rho
sweep}{Graph family G\_\{\textbackslash rho\} and rho sweep}}\label{sec:null-rho}

We define a family of mutation graphs \(G_{\rho}\) parameterized by
\(\rho \in \lbrack 0,1\rbrack\), where \(G_{0} = Q_{6}\) (192 Hamming-1
edges) and \(G_{1} = H(3,4)\) (all 288 single-nucleotide substitution
edges). The weighted mismatch score on \(G_{\rho}\) is

\[F_{\rho(c)} = F_{Hamming-1}(c) + \rho \cdot F_{\text{diagonal}}(c)\]
\protect\phantomsection\label{eq:weighted}{}

where \(F_{\text{diagonal}}\) sums over the 96 within-nucleotide
distance-2 edges. This formulation treats \(Q_{6}\) not as the
biological object but as one endpoint of a continuous interpolation
toward the full mutation graph. We evaluated five values of
\(\rho \in \left\{ 0,0.25,0.5,0.75,1 \right\}\) with \(n =\)10,000
quartet-pattern shuffle null samples per value (common-seed design).
Conservative p-values, computed as \(\frac{k + 1}{n + 1}\) where \(k\)
is the number of null samples scoring below the observed code:
\(p_{\rho = 0} = 0.006\) (\(k = 60\)), \(p_{\rho = 0.25} = 0.0023\)
(\(k = 22\)), \(p_{\rho = 0.5} = 0.0007\) (\(k = 6\)),
\(p_{\rho = 0.75} = 0.0003\) (\(k = 2\)), \(p_{\rho = 1} = 0.0003\)
(\(k = 2\)). The last two are near the Monte Carlo resolution limit
(\(\frac{1}{n + 1} \approx 10^{- 4}\) at \(n =\)10,000); the qualitative
conclusion (that the observed code lies in the extreme low-cost tail at
every \(\rho\)) does not depend on \(n\).

\subsubsection{Table-preserving permutation null (topology
avoidance)}\label{sec:null-topo}

To test whether natural codon reassignment events avoid disrupting amino
acid connectivity, we define a single primary candidate universe of
single-codon reassignments and treat all reasonable alternatives as a
sensitivity analysis. From the standard code \(C\), the primary
candidate set \(\mathcal{M}(C)\) pairs each of the 64 codons with each
of the 20 alternative labels drawn from
\(\mathcal{A}_{20} \cup \left\{ \text{Stop} \right\}\):

\[\mathcal{M}(C) = \left\{ (x,y):x \in \text{ codons},y \in \mathcal{A}_{20} \cup \left\{ \text{Stop} \right\},y \neq C(x) \right\},\quad\left| {\mathcal{M}(C)} \right| = 64 \times 20 = 1280.\]
\protect\phantomsection\label{eq:candidate}{}

This \textbf{21-label, identity-excluded} universe (denoted U1 in
Supplement §S5) gives every codon exactly 20 alternative labels, admits
biologically attested stop-codon reassignment, and excludes identity
moves (\(y = C(x)\)), which contribute no signal. Two strict variants
(the \textbf{amino-acid-only, identity-excluded} universe with
\(\left| \mathcal{M} \right| = 61 \times 19 + 3 \times 20 =\)1,219 (U2)
and the \textbf{stop-inclusive with no-ops} universe with
\(\left| \mathcal{M} \right| = 64 \times 21 =\)1,344 (U4)) are reported
as sensitivity analyses in Supplement §S5; topology-avoidance results
are qualitatively identical under all three.

For each amino acid \(a\), let \(G_{a}^{1}\) denote the subgraph induced
on codons assigned to \(a\) with edges between codons at binary Hamming
distance \(\leq 1\). A reassignment is \emph{topology-breaking} if it
increases the number of connected components of \(G_{a}^{1}\) for any
amino acid \(a\) (the \textbf{increase-in-components} definition
\(\Delta\beta_{0} > 0\) used by the conditional logit; the
\textbf{new-disconnection-in-previously-connected-family} definition is
reported in parallel as a sensitivity check; see Supplement §S3). From
the 27 NCBI translation tables analyzed in this work (Section 2.3.2), we
catalogued the table-specific reassignment events relative to the
standard code; de-duplicating by (codon, target amino acid) tuple yields
the unique-event set on which the hypergeometric and table-preserving
permutation tests operate. The conditional logit model (Section 2.3.5)
operates on the full table-specific event-step list, which preserves
recurrent reassignments across independent lineages.

Two tests were applied: (i) a hypergeometric test treating the observed
events as a sample from the finite candidate landscape (\(N =\)1,280,
\(K = 846\) topology-breaking under the primary cell of \(H(3,4)\)
adjacency with \(\Delta\beta_{0} > 0\) definition, against \(n = 28\)
observed, \(x = 6\) topology-breaking observed); and (ii) a
table-preserving permutation test (\(n =\)10,000, seed 135325) that
permutes target amino acids among a table's reassigned codons,
preserving within-table codon and target structure to address
phylogenetic non-independence. Risk ratios with 95\% log-normal
confidence intervals were computed as \(\text{RR } = (x/n)/(K/N)\). The
full \(2 \times 2\) definition \(\times\) adjacency audit (\(Q_{6}\) vs
\(H(3,4)\) adjacency \(\times\) new-disconnection vs
\(\Delta\beta_{0} > 0\) definitions) is given in Supplement §S3; risk
ratios fall in 0.28--0.33 across all four cells.

\subsubsection{Conditional logit model of reassignment
choice}\label{sec:null-condlogit}

To test whether topology avoidance contributes independent explanatory
power beyond physicochemical optimization, we fit event-level
conditional logit (discrete-choice) models to natural reassignment
events. Each observed reassignment is treated as a choice from the set
of all \(\approx\)1,280 possible single-codon reassignments available at
the current code state. For a table with \(k\) reassignment events, the
code state evolves sequentially: at each step, the model assigns a
probability to each candidate move based on a linear utility score, and
the observed move is compared against all alternatives.

For each candidate move \(m\) from code state \(C\), we computed three
features: (i) \(\Delta_{\text{phys}}\), the change in local Grantham
mismatch cost summed over Hamming-1 edges incident to the reassigned
codon; (ii) \(\Delta_{\text{topo}}\), the total increase in connected
components across all amino acid codon graphs at \(\varepsilon = 1\);
and (iii) \(\Delta_{\text{tRNA}}\), the Hamming distance from the
reassigned codon to the nearest codon already encoding the target amino
acid (a heuristic proxy for tRNA repertoire disruption). The conditional
logit probability of observing move \(m^{\ast}\) is

\[P\left( m^{\ast}~|~\mathcal{N}(C) \right) = \frac{\exp(\mathbf{w}^{\top}\mathbf{x}_{m^{\ast}})}{\sum_{m \in \mathcal{N}(C)}\exp(\mathbf{w}^{\top}\mathbf{x}_{m})}\]
\protect\phantomsection\label{eq:condlogit}{}

where \(\mathbf{w}\) is the weight vector and \(\mathbf{x}_{m}\) is the
feature vector for move \(m\). Features were \(z\)-scored across all
candidates for numerical stability.

Since the temporal ordering of reassignment events within a table is
unknown, we marginalized the likelihood over all \(k!\) orderings for
tables with \(k > 1\) changes:

\[L_{\text{table }} = \frac{1}{k!}\sum_{\sigma \in S_{k}}\prod_{s = 1}^{k}P\left( m_{\sigma(s)}^{\ast}~|~\mathcal{N}(C_{\sigma,s}) \right)\]
\protect\phantomsection\label{eq:orderavg}{}

For tables with \(k \leq 6\) events (including the largest, yeast
mitochondrial, \(k = 6\)), we enumerate all \(k!\) orderings exactly (up
to 720 for \(k = 6\)); for the rare cases of \(k > 6\), we sample 720
random orderings with the seeded RNG. The total likelihood is the
product across all 25 tables with reassignment events (66 total
event-steps).

Four candidate models were compared under \(Q_{6}\) topology;
likelihood-ratio tests were applied only to nested pairs (M1 and M2 are
non-nested alternatives, both nested within M3; M3 is nested in M4): M1
(physicochemistry only, \(w_{\text{topo }} = w_{\text{tRNA }} = 0\)), M2
(topology only, \(w_{\text{phys }} = w_{\text{tRNA }} = 0\)), M3
(physicochemistry + \(Q_{6}\) topology, \(w_{\text{tRNA }} = 0\)), and
M4 (all three features). Two additional verification variants were fit
using the encoding-independent \(H(3,4)\) topology feature in place of
\(Q_{6}\): M2\textsubscript{H(3,4)} (topology only, \(H(3,4)\)) and
M3\textsubscript{H(3,4)} (physicochemistry + \(H(3,4)\) topology). The
\(H(3,4)\) variants test whether M3 dominance is robust to the choice of
topology graph; the
\(\Delta_{topo,H(3,4)\ } = \sum_{a}\left( \beta_{0}\left( G_{a}^{after,H(3,4)} \right) - \beta_{0}\left( G_{a}^{before,H(3,4)} \right) \right)\)
feature replaces the \(Q_{6}\) component-count change. Weights were
estimated by maximum likelihood (scipy.optimize, Nelder--Mead with
L-BFGS-B refinement). Model comparison used the corrected Akaike
Information Criterion (AICc) and likelihood-ratio tests for nested
pairs.

\subsubsection{Fisher--Stouffer test for tRNA
enrichment}\label{sec:null-trna}

For each variant-code organism paired with a phylogenetically proximate
standard-code control, we built a \(2 \times 2\) contingency table
comparing the proportion of tRNA genes encoding the reassigned amino
acid versus all other amino acids. The sampling model treats tRNA gene
counts as draws from a hypergeometric distribution conditional on the
row and column marginals (total tRNAs per organism, total tRNAs for the
focal amino acid), which is appropriate when the question is whether a
specific amino acid's share of the tRNA repertoire differs between two
genomes. Fisher's exact test (one-sided, alternative = ``greater'') was
applied per pairing, and \(p\)-values were combined via Stouffer's \(Z\)
method. To address non-independence from shared control organisms, we
constructed a conflict graph (edges connect pairings sharing an
organism) and enumerated all maximal independent sets (MIS) via the
Bron--Kerbosch algorithm applied to the complement graph. For each MIS
with \(\geq 2\) members we computed the Stouffer combined \(p\)-value
and report the full distribution across MIS (best, median, worst)
together with a pre-specified topology-breaking subset (\(n = 4\)
pairings whose underlying reassignment creates a new amino-acid
disconnection under \(Q_{6}\): yeast mitochondrial Thr,
\emph{Scenedesmus} mitochondrial Leu, \emph{Pachysolen} nuclear Ala,
\emph{Candida} nuclear Ser) as the mechanistic complement.

tRNA gene counts for 18 organisms were obtained by running tRNAscan-SE
2.0.12 \citep{chan2019} with Infernal 1.1.4 on NCBI genome assemblies
(eukaryotic mode for ciliates and yeast; bacterial mode for
\emph{Mycoplasmoides}). The verified set comprises 5 organisms with
translation table 6 (UAA/UAG\(\rightarrow\)Gln; ciliates:
\emph{Tetrahymena thermophila}, \emph{Paramecium tetraurelia},
\emph{Oxytricha trifallax}, \emph{Pseudocohnilembus persalinus},
\emph{Halteria grandinella}), 6 with table 10 (UGA\(\rightarrow\)Cys;
\emph{Euplotes} species), 1 with table 15 (\emph{Blepharisma stoltei};
the genome tested here reads UGA\(\rightarrow\)Trp via a dedicated
suppressor tRNA-Trp(UCA) \citep{singh2023}, despite the legacy NCBI
table-15 UAG\(\rightarrow\)Gln assignment), 1 with table 31 (multiple
stops reassigned; \emph{Blastocrithidia nonstop}), 2 with table 4
(UGA\(\rightarrow\)Trp; \emph{Mycoplasmoides genitalium}, \emph{M.
pneumoniae}), and 3 standard-code controls (\emph{Stentor coeruleus},
\emph{Ichthyophthirius multifiliis}, \emph{Fabrea salina}). One
additional standard-code organism (\emph{Saccharomyces cerevisiae}) was
sourced from GtRNAdb (\url{https://gtrnadb.ucsc.edu}) for context but is
not included in the primary 24-pairing analysis. The dataset comprises
24 pairings across 5 variant genetic codes (NCBI translation tables 4,
6, 10, 15, and 31) and 3 tRNAscan-SE--verified standard-code controls.

\subsection{Retrospective cross-study reanalysis of synthetic recoding
outcomes}\label{sec:methods-codonsafe}

To test whether GF(2)\textsuperscript{6} topology predicts codon
recoding outcomes in synthetic biology experiments, we performed a
retrospective cross-study reanalysis across nine published genome
recoding datasets, with quantitative analysis of eight (the ninth,
\citep{ding2024}, is included for cross-kingdom scope without
quantitative extraction). For each codon substitution reported in these
studies, we computed three topology features under the default encoding
(\(C \mapsto (0,0)\), \(U \mapsto (0,1)\), \(A \mapsto (1,0)\),
\(G \mapsto (1,1)\)):

\begin{enumerate}
\item
  \textbf{Boundary crossing} (\(\varepsilon = 1\)): whether the source
  and target codons lie in different connected components of the amino
  acid's codon graph at Hamming distance 1. Generalized beyond serine to
  all amino acids with potential disconnections (including leucine,
  arginine).
\item
  \textbf{Local mismatch change} (\(\Delta F_{\text{local}}\)): the
  difference in total Grantham distance summed over all Hamming-1
  neighbors between the target and source codon positions:
  \(\Delta F_{\text{local }} = \sum_{n \in N(t)}\Delta(c(t),c(n)) - \sum_{n \in N(s)}\Delta(c(s),c(n))\),
  where \(N(v)\) denotes the 6 Hamming-1 neighbors of vertex \(v\) in
  \(Q_{6}\).
\item
  \textbf{Hamming distance}: the number of bit positions differing
  between source and target vectors in GF(2)\textsuperscript{6}.
\end{enumerate}

Codon positions were extracted from published GenBank genome files by
CDS-level comparison using Biopython \citep{cock2009}, with codon frame
offsets (\texttt{/codon\_start}) handled explicitly. CDS features were
matched between parent and recoded genomes by locus tag. The primary
analysis used the Syn57 dataset \citep{robertson2025syn57}, which
provides a within-study contrast: 37,146 serine recodings cross the
UCN\(\leftrightarrow\)AGY family boundary, while 22,859 alanine
recodings remain within the GCN family. Syn57 RNA-seq differential
expression data (Data file S10) were used as the outcome measure.
Design-to-final genome deviations were identified by comparing the
vS33A7 design genome (Data file S2) against the verified Syn57 genome
(Data file S8), filtering genes with \(>\)10 CDS-level differences as
structural rearrangements (sensitivity analysis across thresholds 3--50
reported in the Supplementary Material). The \citep{ostrov2016} segment
viability data (Table S4) were analyzed as a case-control comparing
segments with and without lethal design exceptions; Bonferroni
correction was applied for three simultaneous tests.

All cross-study reanalysis code is released in the
\texttt{codon\_topo.analysis.codonsafe} subpackage of the
\texttt{codontopo} repository
(\url{https://github.com/biostochastics/codontopo,} version 0.6.1) and
can be run via \texttt{codon-topo\ codonsafe}. Raw data files and
download provenance are documented in
\texttt{data/codonsafe/DATA\_MANIFEST.md}.

\subsection{Structural-preservation index
(visualization-only)}\label{sec:methods-feasibility}

For Figure 5A we use a heuristic \textbf{structural-preservation index}
\(S(m) \in \lbrack 0,1\rbrack\) that combines four discrete indicators
for each candidate single-codon reassignment \(m\): whether the
post-reassignment code preserves the two-fold bit-5 filtration
(\(I_{\text{tw}}\)), whether it preserves the four-fold prefix
filtration (\(I_{\text{ff}}\)), whether Serine remains disconnected at
\(\varepsilon = 1\) (\(I_{\text{ser}}\)), and how many amino acids are
disconnected at \(\varepsilon = 1\) (\(n_{\text{disc}}\)). Explicitly

\[S(m) = 0.25 \cdot I_{\text{tw}(m)} + 0.25 \cdot I_{\text{ff}(m)} + 0.30 \cdot I_{\text{ser}(m)} + w\left( n_{\text{disc}(m)} \right),\]

where \(w(1) = 0.20\), \(w(0) = 0.10\), and \(w(n) = 0.00\) for
\(n \geq 2\). Because every component is discrete, \(S(m)\) takes only
four values across the 1,280-move candidate universe
(\(S \in \left\{ 0.55,0.75,0.80,1.00 \right\}\); 193 / 27 / 738 / 322
events; the four visible bars in Figure 5A). The weights are a
hand-tuned prioritisation of the four structural features and are not
derived from any calibration; the index is descriptive, not inferential,
and is not used in any inferential test in this paper. Exact
implementation: Supplement §S12 and
\texttt{src/codon\_topo/analysis/synbio\_feasibility.py} of the
\texttt{codontopo} repository.

\subsection{Claim hierarchy}\label{sec:claims}

All 15 scientific claims evaluated in this work are registered in a
formal claim hierarchy with pre-specified status categories
(\Cref{tbl:claims}): SUPPORTED (passes rigorous null, \(p < 0.01\)),
EXPLORATORY (hypothesis-generating; trend-level or descriptive),
REJECTED (falsified or pre-rejected in literature), FALSIFIED (directly
contradicted by data), and TAUTOLOGICAL (true by construction). The full
hierarchy with justifications is available in the supplementary
materials.

\begin{longtable}[]{@{}
  >{\raggedright\arraybackslash}p{(\linewidth - 6\tabcolsep) * \real{0.16}}
  >{\raggedright\arraybackslash}p{(\linewidth - 6\tabcolsep) * \real{0.44}}
  >{\centering\arraybackslash}p{(\linewidth - 6\tabcolsep) * \real{0.25}}
  >{\raggedleft\arraybackslash}p{(\linewidth - 6\tabcolsep) * \real{0.15}}@{}}
\caption{Summary of 15 evaluated claims. The implemented null is a
\textbf{quartet-pattern shuffle} (16 first-two-base quartets, each with
its internal AA pattern held atomic and shuffled across quartet slots)
rather than the classical Freeland--Hurst amino-acid permutation across
synonymous codon groups. ``MIS'' = maximal independent set enumeration.
\(p\)-values are the primary summary per claim (MIS median for the tRNA
row; most conservative test elsewhere). Shaded rows are non-inferential:
FALSIFIED and REJECTED rows record claims found to be wrong, and
TAUTOLOGICAL rows record claims that are true by construction and carry
no evidentiary weight.}\tabularnewline
\toprule\noalign{}
\begin{minipage}[b]{\linewidth}\raggedright
\textbf{Status}
\end{minipage} & \begin{minipage}[b]{\linewidth}\raggedright
\textbf{Claim}
\end{minipage} & \begin{minipage}[b]{\linewidth}\centering
\textbf{Null model}
\end{minipage} & \begin{minipage}[b]{\linewidth}\raggedleft
\textbf{\emph{p}-value}
\end{minipage} \\
\midrule\noalign{}
\endfirsthead
\toprule\noalign{}
\begin{minipage}[b]{\linewidth}\raggedright
\textbf{Status}
\end{minipage} & \begin{minipage}[b]{\linewidth}\raggedright
\textbf{Claim}
\end{minipage} & \begin{minipage}[b]{\linewidth}\centering
\textbf{Null model}
\end{minipage} & \begin{minipage}[b]{\linewidth}\raggedleft
\textbf{\emph{p}-value}
\end{minipage} \\
\midrule\noalign{}
\endhead
\bottomrule\noalign{}
\endlastfoot
\multirow{4}{=}{Supported} & Cross-metric coloring optimality &
quartet-pattern shuffle \(\times\) 4 metrics & \(\leq 0.006\) \\
& Per-table optimality (26/27) & Per-table quartet-pattern shuffle +
BH-FDR & \(< 0.05\) \\
& \(\rho\)-robustness (\(\rho \in \lbrack 0,1\rbrack\)) &
Quartet-pattern shuffle (weighted edges) & 0.006 \\
& Topology avoidance (\(Q_{6}\) + \(H(3,4)\)) & Table-pres. perm. +
clade excl. & \(10^{- 4}\) \\
\multirow{5}{=}{Exploratory} & tRNA enrichment (MIS median) &
Fisher--Stouffer (332 MIS) & 0.037 \\
& Bit-position bias (dedup.) & \(\chi^{2}\) uniform & 0.075 \\
& Mechanism boundary conditions & Descriptive & --- \\
& Atchley F3/Serine convergence & Qualitative & --- \\
& Disconnection catalogue (4 cases) & Descriptive & --- \\
Falsified & KRAS--Fano clinical prediction & Fisher--Bonferroni & 1.0 \\
\multirow{3}{=}{Rejected} & Serine min-distance-4 invariant &
Counterexample & --- \\
& PSL(2,7) symmetry & Literature & --- \\
& Holomorphic embedding & Verification & --- \\
\multirow{2}{=}{Tautological} & Two-fold bit-5 filtration & --- & --- \\
& Four-fold prefix filtration & --- & --- \\
\end{longtable}

\protect\phantomsection\label{tbl:claims}{}

\section{Results}\label{sec:results}

\subsection{Cross-metric coloring optimality}\label{sec:res-coloring}

Under the quartet-pattern shuffle null with \(n = 10,000\) permutations,
the standard genetic code is significantly low-cost across all four
physicochemical distance metrics examined (\Cref{tbl:metrics}). For
Grantham distance, the observed score is \(F = 13,477\) versus a null
mean of \(14,954 \pm 628\) (\(z = 2.35\), quantile 0.61\%,
\(p = 0.006\)). The same pattern holds for Miyata distance
(\(F = 235.3\) vs \(275.7 \pm 12.9\), \(z = 3.14\), \(p < 0.001\)),
Woese polar requirement (\(F = 347.3\) vs \(410.5 \pm 23.1\),
\(z = 2.73\), \(p = 0.003\)), and Kyte--Doolittle hydropathy
(\(F = 457.2\) vs \(570.7 \pm 40.9\), \(z = 2.78\), \(p = 0.001\)).
Relative to null expectations, the observed code lowers the mismatch
score by 9.9\%, 14.6\%, 15.4\%, and 19.9\%, respectively. Beta-posterior
credible intervals for the empirical \(p\)-values remain entirely below
0.01 for all four metrics (computed from the conservative
\(\frac{k + 1}{n + 1}\) estimator with \(\text{Beta}(k + 1,n - k + 1)\)
posterior; per-metric intervals are listed in
\texttt{output/coloring\_optimality.json} under
\texttt{beta\_credible\_intervals}). Since \citep{freeland1998}
established optimality using polar requirement alone, the cross-metric
concordance demonstrates that error-minimization is robust across
multiple distinct physicochemical parameterizations, not an artifact of
any single distance metric. Under a weaker degeneracy-only null that
preserves only codon family sizes without maintaining block contiguity,
the signal strengthens substantially (\(z > 9\), \(p < 10^{- 4}\) for
all metrics), confirming that the quartet-pattern shuffle null provides
a stringent test.

\begin{longtable}[]{@{}
  >{\raggedright\arraybackslash}p{(\linewidth - 10\tabcolsep) * \real{0.35}}
  >{\raggedleft\arraybackslash}p{(\linewidth - 10\tabcolsep) * \real{0.11}}
  >{\raggedleft\arraybackslash}p{(\linewidth - 10\tabcolsep) * \real{0.16}}
  >{\raggedleft\arraybackslash}p{(\linewidth - 10\tabcolsep) * \real{0.08}}
  >{\raggedleft\arraybackslash}p{(\linewidth - 10\tabcolsep) * \real{0.15}}
  >{\raggedleft\arraybackslash}p{(\linewidth - 10\tabcolsep) * \real{0.15}}@{}}
\caption{Cross-metric sensitivity analysis. Quartet-pattern shuffle null
(\(n = 10,000\), seed 135325) under four code-independent
physicochemical distance measures and the alignment-derived,
code-dependent ProtSub matrix \citep{jia2021} as a structure-aware
sensitivity analysis. All five measures show significant optimality
after Bonferroni correction at \(\alpha = 0.05\) (per-test threshold
\(\alpha/5 = 0.01\)). Effect size
\(z = \left( \mu_{\text{null }} - F_{\text{obs}} \right)/\sigma_{\text{null}}\).
``Improv.'' = percent improvement of observed score over null mean.
Under a weaker degeneracy-only null (preserving only codon family sizes,
not block structure), all measures yield \(z > 9\) and \(p < 10^{- 4}\).
Under the quartet-pattern shuffle null, ProtSub places the standard code
at the 0.04\% quantile (\(n =\)10,000), the most extreme percentile of
any measure tested.}\tabularnewline
\toprule\noalign{}
\begin{minipage}[b]{\linewidth}\raggedright
\textbf{Metric}
\end{minipage} & \begin{minipage}[b]{\linewidth}\raggedleft
\textbf{Observed}
\end{minipage} & \begin{minipage}[b]{\linewidth}\raggedleft
\textbf{Null mean \(\pm\) SD}
\end{minipage} & \begin{minipage}[b]{\linewidth}\raggedleft
\textbf{\emph{z}}
\end{minipage} & \begin{minipage}[b]{\linewidth}\raggedleft
\textbf{\emph{p}}
\end{minipage} & \begin{minipage}[b]{\linewidth}\raggedleft
\textbf{Improv.}
\end{minipage} \\
\midrule\noalign{}
\endfirsthead
\toprule\noalign{}
\begin{minipage}[b]{\linewidth}\raggedright
\textbf{Metric}
\end{minipage} & \begin{minipage}[b]{\linewidth}\raggedleft
\textbf{Observed}
\end{minipage} & \begin{minipage}[b]{\linewidth}\raggedleft
\textbf{Null mean \(\pm\) SD}
\end{minipage} & \begin{minipage}[b]{\linewidth}\raggedleft
\textbf{\emph{z}}
\end{minipage} & \begin{minipage}[b]{\linewidth}\raggedleft
\textbf{\emph{p}}
\end{minipage} & \begin{minipage}[b]{\linewidth}\raggedleft
\textbf{Improv.}
\end{minipage} \\
\midrule\noalign{}
\endhead
\bottomrule\noalign{}
\endlastfoot
Grantham~\citep{grantham1974} & 13,477 & 14,954 \(\pm\) 628 & 2.35 &
0.006 & 9.9\% \\
Miyata~\citep{miyata1979} & 235.3 & 275.7 \(\pm\) 12.9 & 3.14 &
\textless{} 0.001 & 14.6\% \\
Polar requirement~\citep{woese1973} & 347.3 & 410.5 \(\pm\) 23.1 & 2.73
& 0.003 & 15.4\% \\
Kyte--Doolittle~\citep{kyte1982} & 457.2 & 570.7 \(\pm\) 40.9 & 2.78 &
0.001 & 19.9\% \\
ProtSub~\citep{jia2021}\footnote{Code-dependent (MSA-derived), reported
  as a structure-aware robustness check; see Section 2.2 and
  \citep{digiulio2001}.} & 1089.5 & 1180.2 \(\pm\) 29.7 & 3.05 &
\(4 \times 10^{- 4}\) & 7.7\% \\
\end{longtable}

\protect\phantomsection\label{tbl:metrics}{}

\textbf{Convergent operationalisation across five measures, including
ProtSub.} The four code-independent measures in \Cref{tbl:metrics} are
intercorrelated by construction. Across the 190 unordered amino-acid
pairs (20 standard AAs), pairwise Spearman correlations range from
\(\rho = 0.36\) (Grantham vs Kyte--Doolittle) to \(\rho = 0.73\) (Miyata
vs Woese polar requirement; median \(\rho = 0.49\); full matrix in
Supplement §S20). At the code level (ranking 2,000 random codes drawn
from the quartet-pattern shuffle null by their \(F\)-scores), pairwise
Spearman correlations are higher (range 0.71--0.89, median 0.82), as
expected. To probe whether the panel captures non-identical residual
structure rather than a single latent factor, we computed
\textbf{partial} Spearman correlations controlling simultaneously for
all other measures: five of ten pairs become approximately uncorrelated
after adjustment (\(\left| \rho_{\text{partial}} \right| < 0.1\)),
including Grantham vs polar requirement
(\(\rho_{\text{partial }} = + 0.03\)), polar requirement vs
Kyte--Doolittle (\(+ 0.01\)), and polar requirement vs ProtSub
(\(+ 0.08\)). Substantial shared variance coexists with non-identical
residual structure across the panel; we therefore describe the
cross-metric concordance as a convergent sensitivity envelope rather
than as independent replication. Under the same quartet-pattern shuffle
null, the alignment-derived ProtSub matrix \citep{jia2021} places the
standard code at the most extreme percentile of any measure tested
(\Cref{tbl:metrics}). With five tests at \(\alpha = 0.05\), the
Bonferroni-corrected per-test threshold is 0.01; all five p-values
satisfy it (Grantham \(p = 0.006\); others \(p \leq 0.003\)). ProtSub is
presented as a code-dependent sensitivity analysis rather than as
primary evidence for code-origin optimality (Section 2.2):
\citep{buschmann2026} demonstrate that all substitution-matrix families
(BLOSUM62, AlphaFold-derived AFSM, and 16 others) perform similarly
because they implicitly encode the same physicochemical reality, which
bounds rather than eliminates the \citep{digiulio2001} circularity
concern for alignment-derived matrices.

Decomposing \(F\) by nucleotide position (\Cref{fig:translational},
panel B) reveals that the second codon position contributes 49.3\% of
the total mismatch, the first position 38.2\%, and the wobble position
only 12.5\%. This gradient mirrors the biochemical hierarchy of
mutational impact: second-position changes are the most
physicochemically disruptive, first-position changes are intermediate,
and wobble-position changes are largely synonymous
\citep{woese1965, freeland1998}.

\begin{figure}
\centering
\emfig{Fig1\_core\_topology}
\caption{Core topology of the genetic code in \({\text{GF}(2)}^{6}\).
\textbf{(A)} Persistent homology: connected components (\(\beta_{0}\))
of amino acid codon graphs as a function of Hamming distance threshold
\(\varepsilon\). Serine (bold) is the only amino acid disconnected at
\(\varepsilon = 1\) across all 27 NCBI translation tables
(encoding-invariant); under the default encoding it reconnects at
\(\varepsilon = 4\) (reconnection depth is not encoding-invariant; 16 of
24 encodings give \(\varepsilon = 2\)). \textbf{(B)} Disconnection
catalogue across all translation tables. Each tile denotes an amino acid
disconnected at \(\varepsilon = 1\); the number is the reconnection
\(\varepsilon\) under the default encoding. Serine is universally
disconnected; Leu, Thr, and Ala appear as variant-code disconnections.}
\end{figure}

\protect\phantomsection\label{fig:topology}{}

\begin{figure}
\centering
\emfig{Fig2\_coloring\_optimality}
\caption{Coloring optimality of the genetic code. \textbf{(A)} Null
distribution of Grantham edge-mismatch scores \(F\) under the
quartet-pattern shuffle null (\(n = 10,000\)). The observed standard
code (red line, \(F = 13,477\)) falls at the 0.61th percentile
(\(p = 0.006\)). \textbf{(B)} \(p\)-value across diagonal-edge weight
\(\rho\) from 0 (\(Q_{6}\)) to 1 (full \(H(3,4)\)); all values below
0.01, with optimality strengthening monotonically. The curve
interpolates over 11 sample points shown as markers; the 5 pre-specified
inferential grid values \(\rho \in \left\{ 0,0.25,0.5,0.75,1 \right\}\)
(Table 3) are used for multiplicity correction, and the remaining 6
points are descriptive interpolation shown for shape only. \textbf{(C)}
Per-table quantile for each of 27 NCBI translation tables under their
own quartet-pattern shuffle null (\(n =\)10,000); only table 3 (yeast
mito.) exceeds the 5\% threshold.}
\end{figure}

\protect\phantomsection\label{fig:coloring}{}

\subsection{Robustness across diagonal-edge weight ρ}\label{sec:res-rho}

The weighted score \(F_{\rho}\) (\Cref{eq:weighted}) remains
significantly optimal across all values of \(\rho\) from 0 to 1
(\Cref{tbl:rho}). At \(\rho = 0\) (pure \(Q_{6}\)), \(p = 0.006\)
(\(z = 2.35\)); at \(\rho = 1\) (full \(H(3,4)\), all 288
single-nucleotide edges equally weighted), \(p < 0.001\) (\(z = 3.37\)).
No value of \(\rho\) yields \(p > 0.006\). The optimality signal
strengthens monotonically as diagonal edges are included, addressing the
concern that the \(Q_{6}\) representation ignores approximately
one-third of single-nucleotide mutations.

\begin{longtable}[]{@{}crrrrr@{}}
\caption{Robustness of coloring optimality across diagonal-edge weight
\(\rho\). Quartet-pattern shuffle null, \(n = 10,000\) per \(\rho\). All
\(p\)-values \(< 0.01\). Quantiles reported as empirical rank; values
\(\leq 0.02\%\) indicate \(\leq 2\) null samples below the observed
score.}\tabularnewline
\toprule\noalign{}
\textbf{ρ} & \textbf{Observed} & \textbf{Null mean} & \textbf{Quantile}
& \textbf{\emph{p}-value} & \textbf{\emph{z}} \\
\midrule\noalign{}
\endfirsthead
\toprule\noalign{}
\textbf{ρ} & \textbf{Observed} & \textbf{Null mean} & \textbf{Quantile}
& \textbf{\emph{p}-value} & \textbf{\emph{z}} \\
\midrule\noalign{}
\endhead
\bottomrule\noalign{}
\endlastfoot
\(0\) & 13,477 & 14,954 & 0.61\% & 0.006 & 2.35 \\
\(0.25\) & 15,335 & 17,020 & 0.22\% & 0.002 & 2.77 \\
\(0.5\) & 17,192 & 19,085 & 0.06\% & \textless{} 0.001 & 3.09 \\
\(0.75\) & 19,050 & 21,150 & 0.02\% & \textless{} 0.001 & 3.28 \\
\(1\) & 20,908 & 23,216 & 0.02\% & \textless{} 0.001 & 3.37 \\
\end{longtable}

\protect\phantomsection\label{tbl:rho}{}

\subsection{Per-table optimality preservation}\label{sec:res-pertable}

Of the 27 NCBI translation tables, 26 remain in the top 5\% of their own
quartet-pattern shuffle null after Benjamini--Hochberg FDR correction
(\Cref{fig:coloring}, panel C). The mean quantile across all tables is
1.3\%. The single marginal exception is translation table 3 (yeast
mitochondrial code, 6 codon reassignments; \(p_{\text{BH }} =\)0.073).

The headline number conflates two distinct evidentiary regimes, which we
disaggregate. Variant tables differing from the standard code by
\(\geq 3\) reassignments (``informative-distance'' tables) sit far
enough away that their per-table quartet-shuffle null contains many
distant permutations; tables with \(\leq 2\) reassignments
(``near-standard'' tables) sit so close that their null is dominated by
near-standard permutations. A standard-code-proximity audit (Supplement
§S8) makes the rationale explicit: every variant sits at the extreme
low-\(d_{H}\) tail of its own null (every null draw has
\(d_{H} \geq 30\), while every variant has \(d_{H} \leq 6\)), so no
distance-matched null draw exists. The per-table \(p\)-value therefore
\textbf{cannot} separate proximity-driven from independently-optimized
codes: it measures whether each variant is low-cost \textbf{relative to
its own quartet-shuffle ensemble}, not whether that optimality is
independent of standard-code proximity. The informative-distance vs
near-standard split rests on the absolute \(d_{H} \geq 3\) cutoff, a
descriptive threshold on reassignment distance rather than an identified
boundary between regimes.

Of the 12 informative-distance \textbf{registry entries}, 11 retain
top-5\% placement of their own quartet-shuffle ensembles after BH--FDR
correction, the lone marginal case being yeast mitochondrial. Because
NCBI tables 27 and 28 (Karyorelict and Condylostoma nuclear codes) share
identical sense-codon mappings, this registry corresponds to 11 distinct
sense-codon colorings, 10 of which retain top-5\% placement. Of the 14
near-standard tables, 14 are formally significant under the same
correction, but those results reflect proximity to the standard code
rather than independent variant-specific optimization. The descriptive
claim is therefore that, among variant codes distant enough for the
per-table test to admit a nontrivial range of permutations, all but one
preserve low physicochemical cost relative to their own ensemble. Even
the most reassigned code (table 3, 6 codons reassigned) sits at the
7.4th percentile of its quartet-shuffle null, marginally above the 5\%
threshold but still in the lower tail.

\subsection{Topology avoidance in natural codon
reassignments}\label{sec:res-topo}

We adopt the encoding-independent \(H(3,4)\) Hamming graph as the
primary adjacency for the topology-avoidance test, with the \(Q_{6}\)
subgraph reported as a coordinate-dependent decomposition (motivated
below). Under \(H(3,4)\), of the 1,280 candidate single-codon
relabelings (Methods §2.3.4), 846 (66.1\%) increase the
connected-component count of some amino acid's codon graph relative to
the standard code (\(\Delta\beta_{0} > 0\)). Only 6 of 28 observed
natural reassignment events (21.4\%) do so, yielding a 3.1-fold
depletion (\(\text{RR } = 0.32\), 95\% CI \(\lbrack 0.16,0.66\rbrack\);
hypergeometric \(p =\)1.28 times 10\^{}(-6)\(\); table-preserving
permutation \(p \leq 10^{- 4}\)).

Under \(Q_{6}\) (Hamming-1 adjacency in the default
\({\text{GF}(2)}^{6}\) encoding), the corresponding numbers are 931 of
1,280 candidates (72.7\%) creating new amino-acid disconnections,
against only 6 of 28 observed events (21.4\%); \(Q_{6}\) depletion is
3.4-fold (hypergeometric \(p =\)1.58 times 10\^{}(-8)\(\); permutation
\(p \leq 10^{- 4}\)). The Q\_6 candidate-rate is somewhat higher than
the H(3,4) rate because some \(Q_{6}\)-disconnected pairs become
connected when within-nucleotide Hamming-2 edges are admitted. Both
adjacencies yield highly significant depletion under both
topology-breaking definitions (new disconnection in a previously
connected family; \(\Delta\beta_{0} > 0\) increase in components), with
risk ratios 0.28--0.33 across the four cells of the \(2 \times 2\)
definition \(\times\) adjacency audit (Supplement §S3) --- the avoidance
is not specific to either the hypercube subgraph or either definition.
We give precedence to the \(H(3,4)\) result because \(Q_{6}\) is
encoding-dependent: across all 24 base-to-bit bijections, the \(Q_{6}\)
candidate-landscape rate ranges from 36\% (8 encodings) to 73\% (default
encoding), with median hypergeometric \(p = 6 \times 10^{- 6}\) but
\(p > 0.5\) in 8 of 24 encodings (Supplement §S4). \(H(3,4)\), by
contrast, depends only on nucleotide identity and is robust by
construction.

Following \citep{sengupta2007}, we performed clade-exclusion sensitivity
analysis, iteratively removing each major taxonomic group (all ciliates,
all metazoan mitochondria, all CUG-clade yeasts, etc.) and retesting.
The depletion is highly significant in every regime (\(p < 10^{- 5}\)),
but requires cautious interpretation: the topology-breaking events are
not broadly distributed across clades. The full 7-row clade-exclusion
table is persisted for both the \(Q_{6}\) / new-disconnection adjacency
and the primary \(H(3,4)\) / \(\Delta\beta_{0} > 0\) adjacency
(Supplement §S9). For the yeast-mitochondrial-excluded regime the direct
calculation is: 4 of the 6 \(H(3,4)\) breakers come from a single clade
(yeast mitochondrial, NCBI translation table 3; all 4 of its
de-duplicated reassignments are \(H(3,4)\)-breakers), and excluding it
leaves \(\frac{2}{24} = 8.3\)\% breakers, moving the \(H(3,4)\),
\(\Delta\beta_{0} > 0\) hypergeometric \(p\) from \(\)1.28 times
10\^{}(-6)\(\) to \(\approx 4.2 \times 10^{- 9}\) (direct one-sided
calculation on \(N = 1280\), \(K = 846\), \(n = 24\), \(x = 2\)). The
\(Q_{6}\) / new-disconnection variant of the same single-row calculation
gives \(p \approx 3.6 \times 10^{- 11}\) (Supplement §S9) and is not the
primary cell. The clade-exclusion robustness is thus a denominator
effect: removing clades that contributed zero or one breaker each leaves
the breakage rate essentially unchanged, rather than evidence of
repeated, independently arising topology-breaking events across many
lineages. The cross-lineage signal is therefore the \textbf{avoidance}
of topology-breaking moves, supported by 22 of 28 events across many
lineages; the cross-clade structure of the breaker set itself is
underpowered.

NCBI translation table 3 (yeast mitochondrial) appears as the marginal
case in three diagnostics: per-table optimality (§3.3), the
lowest-ranking observed move under the M3 conditional logit (§3.5), and
the topology-breaker concentration here. These three observations share
a \(k\)-magnitude effect by construction (table 3 has the most
reassignments in the registry, and is the only variant code requiring
acquisition of a novel tRNA from a different parent:
tRNA\textsuperscript{Thr} derived from tRNA\textsuperscript{His}
\citep{su2011}), so we do not interpret their convergence as independent
confirmation of any specific selectionist mechanism. The pattern is
consistent with topology-breaking reassignments being feasible but
appearing concentrated in unusually perturbed translation systems.

All four cells of the adjacency \(\times\) definition audit
(\Cref{tbl:topo-audit}) yield hypergeometric \(p < 10^{- 5}\) with risk
ratios 0.28--0.33, so the topology-breaking definition matters little.
The encoding-sensitivity audit matters more: the \(H(3,4)\) result is
constant across all 24 base-to-bit bijections, but under the 8 encodings
placing the \(Q_{6}\) candidate rate near 36\% the observed rate
(21--36\%) does not differ from it (\(p > 0.5\)). We therefore treat
\(H(3,4)\) as the primary topology-avoidance result and \(Q_{6}\) as a
coordinate-dependent decomposition (Supplement §S4). Denominator
sensitivity (alternative candidate universes \(|M| = 1219\),
\(|M| = 1344\)) is reported in Supplement §S5 and does not change the
qualitative conclusion.

\begin{longtable}[]{@{}lrrr@{}}
\caption{Topology avoidance in natural codon reassignments under the
encoding-independent \(H(3,4)\) Hamming graph (primary cell).
Topology-breaking = \(\Delta\beta_{0} > 0\), an increase in the total
number of connected components summed across amino acid codon graphs.
The \textbf{Expected (chance)} column shows what would be observed under
the null hypothesis that the 28 observed events are drawn uniformly at
random from the 1,280-move candidate landscape:
\(E\left\lbrack \text{breaks} \right\rbrack = 28 \times K/N \approx 18.5\)
versus 6 observed. RR = risk ratio (observed rate / candidate rate) with
95\% log-normal CI; depletion fold = (candidate rate)/(observed rate).
The remaining three cells of the \(2 \times 2\) definition \(\times\)
adjacency audit (\(Q_{6}\) vs \(H(3,4)\) \(\times\) new-disconnection vs
\(\Delta\beta_{0} > 0\)) are reported in Supplement §S3 (Table 5 of this
manuscript reproduces them inline), with risk ratios in the range
0.28--0.33 and hypergeometric \(p < 10^{- 5}\)
throughout.}\tabularnewline
\toprule\noalign{}
\textbf{Quantity} & \textbf{Observed} & \textbf{Candidate landscape} &
\textbf{Expected (chance)} \\
\midrule\noalign{}
\endfirsthead
\toprule\noalign{}
\textbf{Quantity} & \textbf{Observed} & \textbf{Candidate landscape} &
\textbf{Expected (chance)} \\
\midrule\noalign{}
\endhead
\bottomrule\noalign{}
\endlastfoot
Topology-breaking (\(\Delta\beta_{0} > 0\)) & 6 (21.4\%) & 846 (66.1\%)
& 18.5 (66.1\%) \\
Topology-preserving (\(\Delta\beta_{0} = 0\)) & 22 (78.6\%) & 434
(33.9\%) & 9.5 (33.9\%) \\
Total & 28 & 1,280 & 28 \\
Depletion fold (cand. rate / obs. rate) & --- & 3.1\(\times\) &
1.0\(\times\) \\
RR (95\% CI) & --- & 0.32 (0.16--0.66) & 1.00 (ref.) \\
Hypergeometric \emph{p} (one-sided) & --- & \(1.28 \times 10^{- 6}\) &
(null reference) \\
Permutation \emph{p} (table-pres.) & --- & \(\leq 10^{- 4}\) & (null
reference) \\
\end{longtable}

\protect\phantomsection\label{tbl:topo}{}

\begin{longtable}[]{@{}llrrr@{}}
\caption{Sensitivity analysis: the \(2 \times 2\) topology-breaking
definition \(\times\) adjacency audit. All four cells share the same
denominator (1,280 candidates) and the same observed-event total (28
de-duplicated reassignments); the observed topology-breaking count is
5--7 of 28 across cells. All four cells yield depletion in the same
direction with comparable risk ratios (0.28--0.33). The \(H(3,4)\)
Hamming graph is encoding-independent; the \(Q_{6}\) subgraph is
encoding-dependent and 8 of 24 base-to-bit bijections give no \(Q_{6}\)
depletion at all (Supplement §S4). The main text uses the \(H(3,4)\),
\(\Delta\beta_{0} > 0\) cell as primary.}\tabularnewline
\toprule\noalign{}
\textbf{Adjacency} & \textbf{Definition} & \textbf{\(K/N\)} &
\textbf{Hyper. \emph{p}} & \textbf{RR (95\% CI)} \\
\midrule\noalign{}
\endfirsthead
\toprule\noalign{}
\textbf{Adjacency} & \textbf{Definition} & \textbf{\(K/N\)} &
\textbf{Hyper. \emph{p}} & \textbf{RR (95\% CI)} \\
\midrule\noalign{}
\endhead
\bottomrule\noalign{}
\endlastfoot
\(H(3,4)\) (primary) & \(\Delta\beta_{0} > 0\) & 846 / 1,280 &
\(1.3 \times 10^{- 6}\) & 0.32 (0.16--0.66) \\
\(H(3,4)\) & new disconnection & 822 / 1,280 & \(5.0 \times 10^{- 7}\) &
0.28 (0.13--0.62) \\
\(Q_{6}\) & \(\Delta\beta_{0} > 0\) & 963 / 1,280 &
\(2.2 \times 10^{- 8}\) & 0.33 (0.17--0.63) \\
\(Q_{6}\) & new disconnection & 931 / 1,280 & \(1.6 \times 10^{- 8}\) &
0.29 (0.14--0.60) \\
\end{longtable}

\protect\phantomsection\label{tbl:topo-audit}{}

\begin{figure}
\centering
\emfig{Fig3\_evolutionary\_evidence}
\caption{Evolutionary evidence for structural constraints on codon
reassignment. \textbf{(A)} Bit-position bias: distribution of bit-flips
across \({\text{GF}(2)}^{6}\) coordinates in natural reassignment
events; dashed line = uniform expectation. \textbf{(B)} Evolutionary
depth calibration: reconnection \(\varepsilon\) vs estimated divergence
age (log scale) --- 6 calibration points spanning 4 amino acids
(universal Serine as the pre-LUCA baseline plus variant-code Thr, Leu,
and Ala; Serine and Leucine each contribute two points from different
lineages). \textbf{(C)} Topology avoidance shown for the
encoding-independent \(H(3,4)\) Hamming graph (left panel, primary,
using the \textbf{increase-in-components} \(\Delta\beta_{0} > 0\)
definition matched to the conditional-logit feature) and the
encoding-dependent \(Q_{6}\) subgraph (right panel, sensitivity, using
the \textbf{new-disconnection} definition for continuity with the
paper's \(Q_{6}\) narrative). Under \(H(3,4)\), observed natural
reassignments break topology at 21.4\% versus 66.1\% of the candidate
landscape (RR 0.32, permutation \(p \leq 10^{- 4}\)). The four-cell
\(2 \times 2\) audit (definition \(\times\) adjacency) is reported in
Supplement §S3 and yields risk ratios 0.28--0.33 across all cells.
\textbf{(D)} tRNA enrichment: rank of the reassigned amino acid among
all 20 AAs by tRNA gene proportion in variant-code vs standard-code
organism pairings; rank 1 = most enriched.}
\end{figure}

\protect\phantomsection\label{fig:topo}{}

\subsection{Explanatory modeling: topology as an independent
predictor}\label{sec:res-condlogit}

To test whether topology avoidance is reducible to physicochemical
optimization, we fit the six event-level conditional logit models of
\Cref{sec:null-condlogit} to the 66 reassignment event-steps across 25
variant-code tables, treating each observed reassignment as a choice
among \(\approx\)1,280 candidate single-codon moves (four share the
encoding-dependent \(Q_{6}\) topology feature; two are \(H(3,4)\)
verification variants). The combined physicochemistry-plus-topology
model under \(Q_{6}\) topology (M3) was strongly favored over all
alternatives by AICc (\(\text{AICc } = 721.4\)), outperforming the
physicochemistry-only model (M1; \(\Delta\text{AICc } = 108.2\)) and the
topology-only model (M2; \(\Delta\text{AICc } = 89.1\)). The
encoding-independent verification model M3\textsubscript{H(3,4)} also
outperforms the physicochemistry-only baseline
(\(\Delta\text{AICc}\left( M1\  \rightarrow \ M3{H(3,4)} \right) = 91.3\))
and the H(3,4)-topology-only baseline
(\(\Delta\text{AICc}\left( M2{H(3,4)} \rightarrow \ M3{H(3,4)} \right) = 95.1\)),
so the conditional-logit thesis does not depend on the choice of
topology graph. Adding a heuristic tRNA-complexity proxy did not improve
fit (M3\(\rightarrow\)M4: \(\Delta\text{AICc } = 2.1\)).

\begin{longtable}[]{@{}
  >{\centering\arraybackslash}p{(\linewidth - 10\tabcolsep) * \real{0.1667}}
  >{\raggedright\arraybackslash}p{(\linewidth - 10\tabcolsep) * \real{0.1667}}
  >{\centering\arraybackslash}p{(\linewidth - 10\tabcolsep) * \real{0.1667}}
  >{\centering\arraybackslash}p{(\linewidth - 10\tabcolsep) * \real{0.1667}}
  >{\centering\arraybackslash}p{(\linewidth - 10\tabcolsep) * \real{0.1667}}
  >{\centering\arraybackslash}p{(\linewidth - 10\tabcolsep) * \real{0.1667}}@{}}
\caption{Event-level conditional logit model comparison. Six models
share the same candidate set (\(\approx\)1,280 single-codon moves per
event-step) and the same likelihood structure marginalised over event
orderings; they differ only in the feature subset used. The first four
models use the encoding-dependent \(Q_{6}\) topology feature (the legacy
primary, paired with the encoding sweep in §S4); the two M-variants
ending in \(H(3,4)\) replace \(\Delta_{topo,Q\_ 6}\) with the
encoding-independent \(\Delta_{topo,H(3,4)}\) feature, providing the
encoding-robustness verification. M3 (phys + topo) is favoured under
both topology graphs; the gap between M3\textsubscript{Q\_6} and
M3\textsubscript{H(3,4)} reflects which topology feature carries more
information per parameter, not whether topology is itself meaningful.
\(k\) = number of estimated parameters. Bold = best model by
AICc.}\tabularnewline
\toprule\noalign{}
\begin{minipage}[b]{\linewidth}\centering
Model
\end{minipage} & \begin{minipage}[b]{\linewidth}\raggedright
Description
\end{minipage} & \begin{minipage}[b]{\linewidth}\centering
\(k\)
\end{minipage} & \begin{minipage}[b]{\linewidth}\centering
Log \(\mathcal{L}\)
\end{minipage} & \begin{minipage}[b]{\linewidth}\centering
AICc
\end{minipage} & \begin{minipage}[b]{\linewidth}\centering
\(\Delta\)AICc
\end{minipage} \\
\midrule\noalign{}
\endfirsthead
\toprule\noalign{}
\begin{minipage}[b]{\linewidth}\centering
Model
\end{minipage} & \begin{minipage}[b]{\linewidth}\raggedright
Description
\end{minipage} & \begin{minipage}[b]{\linewidth}\centering
\(k\)
\end{minipage} & \begin{minipage}[b]{\linewidth}\centering
Log \(\mathcal{L}\)
\end{minipage} & \begin{minipage}[b]{\linewidth}\centering
AICc
\end{minipage} & \begin{minipage}[b]{\linewidth}\centering
\(\Delta\)AICc
\end{minipage} \\
\midrule\noalign{}
\endhead
\bottomrule\noalign{}
\endlastfoot
\textbf{M3} & Phys + topo (Q\_6, primary) & 2 & \(-358.6\) &
\textbf{721.4} & \textbf{0.0} \\
M4 & Phys + topo (Q\_6) + tRNA proxy & 3 & \(-358.6\) & 723.5 & 2.1 \\
M3 & Phys + topo (H(3,4) verification) & 2 & \(-367.1\) & 738.4 &
16.9 \\
M2 & Topo only (Q\_6) & 1 & \(-404.2\) & 810.5 & 89.1 \\
M1 & Phys only & 1 & \(-413.8\) & 829.7 & 108.2 \\
M2 & Topo only (H(3,4)) & 1 & \(-415.7\) & 833.5 & 112.1 \\
\end{longtable}

\protect\phantomsection\label{tbl:condlogit}{}

Likelihood-ratio tests confirmed that each feature class adds
substantial explanatory value to the other: adding topology to
physicochemistry yields LR \(= 110.4\) (\(p \ll 10^{- 10}\)), and adding
physicochemistry to topology yields LR \(= 91.2\) (\(p \ll 10^{- 10}\)).
The two feature classes are only weakly associated across the full
candidate landscape, indicating limited confounding.

In the best-fitting model (M3), observed natural reassignments
preferentially populate moves that reduce local physicochemical mismatch
(\({\hat{\beta}}_{\text{phys }} = - 0.004\) per Grantham unit) and
strongly avoid moves that increase amino acid codon-family disconnection
(\({\hat{\beta}}_{\text{topo }} = - 3.32\) per additional connected
component, with the conditional-logit feature evaluated under \(Q_{6}\)
adjacency using the \textbf{increase-in-components}
(\(\Delta\beta_{0} > 0\)) convention). Because \(Q_{6}\) adjacency is
encoding-dependent (8 of 24 base-to-bit bijections give no \(Q_{6}\)
topology depletion at the landscape level; \Cref{sec:res-topo}), the
encoding-independent \(H(3,4)\) refit is the key check. The
encoding-robustness comparison gives
\(\Delta\text{AICc}\left( M1\  \rightarrow \ M3{H(3,4)} \right) = 91.3\)
and
\(\Delta\text{AICc}\left( M2{H(3,4)} \rightarrow \ M3{H(3,4)} \right) = 95.1\),
both well above the conventional ΔAICc \textgreater{} 10 threshold
treated as strong evidence in the model-comparison literature
\citep{burnham2002}, and similar in magnitude to the \(Q_{6}\)
counterparts
(\(\Delta\text{AICc}(M1\  \rightarrow \ M3Q\_ 6) = 108.2\)). This
threshold was empirically calibrated on linear-regression
model-comparison contexts; we use it here as a conventional reference
rather than as a formally calibrated cut-off in the conditional-logit
setting. The M3 dominance is therefore not an artifact of the \(Q_{6}\)
encoding choice. Under M3, observed natural reassignments rank on
average at the 89.5th percentile among all candidate moves
(\Cref{fig:condlogit}, panel B), with the recurrent
UGA\(\rightarrow\)Trp reassignment consistently ranking above the 98th
percentile. The one outlier is the yeast mitochondrial
CUU\(\rightarrow\)Thr reassignment (30th percentile), consistent with
NCBI translation table 3 being the sole marginal exception in the
per-table optimality analysis (\Cref{sec:res-pertable}). The
non-significance of the heuristic tRNA-distance proxy (LR = 0.12,
\(p = 0.73\)) speaks to the inadequacy of
Hamming-distance-to-nearest-target-AA-codon as a proxy for tRNA-mediated
mechanistic feasibility; it should not be interpreted as evidence
against tRNA effects in general (see \Cref{sec:res-trna} for direct
tRNA-gene-count tests). A parametric predictive simulation under M3
reproduced the observed topology-breaking rate at the event-step level
under \(Q_{6}\) adjacency using the \textbf{increase-in-components}
(\(\Delta\beta_{0} > 0\)) definition (observed 5/66 event-steps = 0.076
vs simulated mean 0.077 \(\pm\) 0.033 over \(n =\)10,000 draws;
parametric predictive \(p = 0.60\); Supplement §S19.5, Figure S8 panel
C). This event-step / \(Q_{6}\) / \(\Delta\beta_{0} > 0\) rate differs
from the lineage-collapsed 6/28 = 0.21 rate quoted for \(H(3,4)\)
adjacency with the \textbf{increase-in-components}
(\(\Delta\beta_{0} > 0\)) definition in §3.4 (\Cref{sec:res-topo}) ---
the same definition, but a different denominator (66 event-steps vs 28
events) and a different adjacency (\(Q_{6}\) vs \(H(3,4)\)). The
predictive check supports model calibration on its own
denominator/adjacency/definition combination rather than in-sample AICc
improvement alone.

Like all conditional logit models, M1--M4 assume Independence of
Irrelevant Alternatives (IIA): the relative probability between any two
candidate moves is unaffected by adding or removing other candidates. We
use the model as an explanatory rather than predictive tool: the
reported ΔAICc and likelihood-ratio statistics test whether topology
adds explanatory value beyond physicochemical cost, not whether the
model accurately predicts which specific reassignment will occur next;
full IIA discussion appears in Supplement §S6. A separate concern is
that the unrestricted candidate set (\(\approx\)1,280 single-codon
moves) admits strongly deleterious alternatives (reassigning AUG-Met,
multi-codon changes implicit in single-step framing, reassignments to
stop in essential codons) that selection has already removed. Models
with strongly negative coefficients on \(\Delta_{\text{topo}}\) and
\(\Delta_{\text{phys}}\) may therefore partly rediscover that natural
reassignments are not catastrophic, inflating apparent ΔAICc gaps. To
bound this concern we refit M1--M4 (and the \(H(3,4)\) verification
variants) on a \textbf{restricted candidate set} containing only
candidates whose target amino acid is already accessible at Hamming
distance \(\leq d\) from the reassigned codon
(\(\Delta_{\text{tRNA }} \leq d\), with \(d \in \left\{ 1,2,3 \right\}\)
reported as a \textbf{bracketing} sensitivity rather than a designated
primary; Supplement §S6.1). The strictest cut, \(d = 1\), is the most
defensible biological-plausibility floor (canonical wobble tolerance);
\(d = 3\) is a loose upper bound. Under \(d = 2\) the candidate set
shrinks from \(\approx\)1,280 to \(\approx 727\) candidates per choice
set, and M3 still dominates with
\(\Delta\text{AICc}(M1\  \rightarrow \ M3) = 60\) and
\(\Delta\text{AICc}(M2\  \rightarrow \ M3) = 77\), both well above the
conventional ΔAICc\textgreater10 reference. The qualitative explanatory
claim survives even the most stringent \(d = 1\) filter (\(\approx 275\)
candidates per choice set), where
\(\Delta\text{AICc}(M1\  \rightarrow \ M3) = 13.8\) shrinks but remains
above the conventional 10 threshold and
\(\Delta\text{AICc}(M2\  \rightarrow \ M3) = 73\) remains large; under
the looser \(d = 3\) filter (\(\approx 1088\) candidates), the gap
recovers to \(\Delta\text{AICc}(M1\  \rightarrow \ M3) = 95\). As
expected, tighter cuts shrink the ΔAICc gaps, but the qualitative claim
``topology adds explanatory value beyond physicochemistry'' is robust at
every threshold. We therefore read the unrestricted-set ΔAICc magnitudes
as \textbf{upper bounds} on the topology--physicochemistry separation
and treat the \(d = 1\) floor as the most defensible effect-size
bracket. (We do not interpret
\(\Delta\text{AICc}(M3\  \rightarrow \ M4)\) from the restricted set:
because the restriction is itself defined on \(\Delta_{\text{tRNA}}\),
the M4 tRNA-feature distribution shifts mechanically with the filter, so
the M4 comparison is informative only under the unrestricted set, where
it remains uninformative.)

Conditional-logit clade-exclusion sensitivity (refitting M1--M4 with
each major clade dropped, matching the regime applied to the
topology-avoidance hypergeometric in \citep{sengupta2007}) is reported
in Supplement §S7. Across all seven exclusion regimes, M3 is
consistently favored over M1 (\(\Delta\text{AICc } \geq 49\), median
\(101\), max \(117\)) and over M2 (\(\Delta\text{AICc } \geq 32\),
median \(88\)). Even the minimum
\(\Delta\text{AICc}\)(M1\(\rightarrow\)M3) of 49 (from excluding all
metazoan mitochondrial codes, leaving 40 events) sits comfortably above
the conventional ΔAICc\textgreater10 reference threshold; the
qualitative conclusion (that topology adds explanatory value beyond
physicochemistry) survives every clade-exclusion regime. We rely most
heavily on the parametric predictive calibration check reported above
(\(p = 0.60\)), since it does not require importing a threshold from a
different statistical regime.

\begin{figure}
\centering
\emfig{Fig5\_condlogit}
\caption{Event-level explanatory modeling of natural codon
reassignments. \textbf{(A)} AICc comparison of the six candidate
conditional-logit models (M1--M4 under \(Q_{6}\) plus the two
\(H(3,4)\)-verification variants; likelihood-ratio tests were applied
only to nested pairs); lower is better. The combined
physicochemistry-plus-topology model M3 has the lowest AICc; the
single-feature baselines M1 and M2 are 89.1--112.1 AICc units worse, the
\(H(3,4)\)-topology verification model M3\textsubscript{H(3,4)} is 16.9
units worse than M3\textsubscript{Q\_6}, and M4 (adding the heuristic
tRNA-distance proxy) is 2.1 units worse than M3. \textbf{(B)}
Distribution of observed move percentile ranks under M1
(physicochemistry only, light) versus M3 (combined, dark); M3
concentrates ranks toward the top of the candidate set. Dashed line =
chance expectation (50th percentile). \textbf{(C)} Likelihood-ratio
tests for each feature class added to its complement; both topology and
physicochemistry contribute highly significant independent information
(\(p < 0.001\)); the heuristic tRNA proxy does not (\(p = 0.73\)).}
\end{figure}

\protect\phantomsection\label{fig:condlogit}{}

\subsection{tRNA enrichment for reassigned amino
acids}\label{sec:res-trna}

Organisms with variant genetic codes tend to show elevated tRNA gene
copy numbers for the reassigned amino acid relative to standard-code
controls. Across 24 variant--control pairings derived from 24 organisms
with mixed provenance (15 of 18 tRNAscan-SE--verified genomes plus 9
organism-slots from literature/GtRNAdb/annotation; the three
tRNAscan-verified genomes not in this analysis are boundary-case
mechanisms discussed in §4.3; full 24-row catalogue with per-pair 2\ensuremath{\times}2
counts in Supplement §S10.1), Fisher's exact test combined via
Stouffer's \(Z\) method yields \(p =\)1.31 times 10\^{}(-7)\(\)
(\(Z = 5.15\)) for the all-pairings analysis; this pooled value does not
adjust for shared control organisms. To address the non-independence
directly, we enumerated all maximal independent sets (MIS) of the
conflict graph via Bron--Kerbosch on the complement, obtaining 332 MIS
of size 6. Across the 332 sets the Stouffer combined \(p\)-value has
median 0.037, best 0.012, and worst 0.104; 264 of 332 (79.5\%) sets fall
below the 0.05 threshold and 0 of 332 below 0.01
(\Cref{tbl:trna-tests}). A pre-specified topology-breaking subset (n=4:
yeast mitochondrial Thr, \emph{Scenedesmus} mitochondrial Leu,
\emph{Pachysolen} nuclear Ala, \emph{Candida} nuclear Ser) is null
(Stouffer \(p = 0.387\)), so the direct mechanistic link between
tRNA-gene duplication and codon-family disconnection is not supported by
the topology-breaking subset in isolation, and the elevated-tRNA pattern
reflects a heterogeneous accommodation-of-decoding signal across
variant-code lineages rather than a specific
compensation-for-disconnection mechanism. The pattern is also not
universal at the mechanistic level: \emph{Blastocrithidia nonstop} (NCBI
translation table 31) achieved UGA\(\rightarrow\)Trp via anticodon stem
shortening rather than gene duplication \citep{kachale2023}, and
\emph{Mycoplasmoides} species with UGA\(\rightarrow\)Trp use a single
tRNA-Trp with anticodon modification. These boundary cases motivate the
genome-size-dependent mechanistic landscape discussed in
\Cref{sec:disc-trna}; the full 18-organism tRNAscan catalogue
(including the three boundary-case genomes not in the pairings) and
per-lineage mechanistic detail (\emph{Tetrahymena thermophila} with 54
Gln tRNAs including 39 first-pass suppressors; six \emph{Euplotes}
species carrying TCA-anticodon tRNA-Cys reading UGA; and the
\emph{Blastocrithidia}/\emph{Mycoplasmoides} boundary cases) are given
in Supplement §S10, with the 24-row pairing/provenance table in §S10.1.

\begin{longtable}[]{@{}lccr@{}}
\caption{Summary of tRNA enrichment tests across 24 pairings among 24
organisms with mixed provenance (15 of 18 tRNAscan-SE-verified genomes +
9 organism-slots from literature/GtRNAdb/annotation; three
tRNAscan-verified genomes not in this analysis --- \emph{Blastocrithidia
nonstop}, \emph{Mycoplasmoides genitalium}, \emph{Mycoplasmoides
pneumoniae} --- are boundary-case mechanisms discussed in §4.3). MIS =
maximal independent set, enumerated via the Bron--Kerbosch algorithm
applied to the complement of the conflict graph (edges = shared
organisms). The 24-pairing graph admits 332 MIS, all of size 6. Across
the 332 sets the Stouffer combined \(p\)-value distribution has median
0.037, best 0.012, worst 0.104; 264 of 332 (79.5\%) fall below the 0.05
threshold and 0 of 332 below 0.01. The topology-breaking subset
restricts to the four pairings whose underlying reassignment creates a
new amino-acid disconnection under \(Q_{6}\) (yeast mitochondrial Thr,
\emph{Scenedesmus} mitochondrial Leu, \emph{Pachysolen} nuclear Ala,
\emph{Candida} nuclear Ser); this subset yields Stouffer \(p = 0.387\)
and is null, so the tRNA-duplication mechanism is not directly evidenced
by the topology-breaking cases in isolation. The all-pairings
(\(n = 24\)) Stouffer \(p\) is not adjusted for shared control organisms
and is retained as the pooled point estimate only.}\tabularnewline
\toprule\noalign{}
\textbf{Test} & \textbf{\(n\)} & \textbf{Statistic} &
\textbf{\emph{p}-value} \\
\midrule\noalign{}
\endfirsthead
\toprule\noalign{}
\textbf{Test} & \textbf{\(n\)} & \textbf{Statistic} &
\textbf{\emph{p}-value} \\
\midrule\noalign{}
\endhead
\bottomrule\noalign{}
\endlastfoot
Fisher+Stouffer (all pairings) & 24 & \(Z = 5.10\) &
\(1.7 \times 10^{- 7}\) \\
MIS median (across all 332 MIS) & 6 & \(Z \approx 1.79\) & 0.037 \\
\textbf{MIS worst-case} & \textbf{6} & \(\mathbf{Z \approx 1.26}\) &
\textbf{0.104} \\
MIS best-case & 6 & \(Z \approx 2.14\) & 0.016 \\
Topology-breaking subset & 4 & \(Z \approx 0.16\) & 0.435 \\
\end{longtable}

\protect\phantomsection\label{tbl:trna-tests}{}

\subsection{Exploratory observations}\label{sec:res-exploratory}

Three exploratory observations are recorded here without inferential
weight (full detail in Supplement §S23): (i) an apparent bit-position
bias in codon reassignments (\(\chi^{2}p = 0.006\) under a uniform null)
attenuates after de-duplication (\(p = 0.075\)) and vanishes under
codon-preserving permutation (\(p = 1.0\)), indicating that the signal
reflects which codons are recurrently reassigned rather than a
positional preference; (ii) a systematic catalogue across all 27 NCBI
translation tables identifies four variant-code amino-acid
disconnections at \(\varepsilon = 1\) beyond the universal Serine
disconnection: Thr in yeast mitochondrial (table 3), Leu in
chlorophycean mitochondrial (tables 16, 22), Ala in \emph{Pachysolen
tannophilus} (table 26), and a tripartite Ser in \emph{Candida}-clade
nuclear (table 12); and (iii) Serine has the most extreme Atchley Factor
3 score (\(F_{3} = - 4.760\)) and is uniquely disconnected under every
encoding, but these two views of Serine's anomaly are not fully
independent, since both reflect its disproportionate codon diversity
relative to its physicochemical footprint.

\subsection{Retrospective cross-study reanalysis of synthetic genome
recoding outcomes}\label{sec:res-codonsafe}

To assess whether the structural properties identified in Sections
3.1--3.5 translate to measurable phenotypic consequences in synthetic
biology, we performed a retrospective cross-study reanalysis of nine
published genome recoding datasets, with quantitative analysis of eight
(\(>\)217,000 codon-level observations; \Cref{tbl:codonsafe-datasets}).
Each codon substitution was classified by its GF(2)\textsuperscript{6}
topology: whether it crosses a connected-component boundary at
\(\varepsilon = 1\), the change in local physicochemical mismatch cost
(\(\Delta F_{\text{local}}\), Grantham distance), and the Hamming
distance between source and target vectors.

\begin{longtable}[]{@{}p{3.5cm}p{3.6cm}>{\raggedleft\arraybackslash}p{1.2cm}>{\centering\arraybackslash}p{2.5cm}p{3.5cm}@{}}
\caption{Published genome recoding datasets analyzed. ``Boundary?''
indicates whether the synonymous codon swaps cross a connected-component
boundary at \(\varepsilon = 1\) in GF(2)\textsuperscript{6}.
\textsuperscript{†}bioRxiv preprint; not peer-reviewed at time of
writing. Among quantitatively analyzed datasets, Syn57 provides a
\textbf{topology-fixed exposure} (every Ser swap crosses the
UCN\(\leftrightarrow\)AGY boundary; no Ala swap does) paired with a null
gene-level transcriptomic test; the codon-level
\(\Delta F_{\text{local}}\) contrast within Syn57 is design-confounded
and reported descriptively (see §3.8.1.2). \citep{ding2024} is included
qualitatively for cross-kingdom scope without quantitative extraction.
The Syn57 row total of 60,240 codon-level observations comprises 37,146
serine recodings + 22,859 alanine recodings + 235 stop-codon
recodings.}\tabularnewline
\toprule\noalign{}
\textbf{Study} & \textbf{Swap type} & \textbf{\(n\)} &
\textbf{Boundary?} & \textbf{Key result} \\
\midrule\noalign{}
\endfirsthead
\toprule\noalign{}
\textbf{Study} & \textbf{Swap type} & \textbf{\(n\)} &
\textbf{Boundary?} & \textbf{Key result} \\
\midrule\noalign{}
\endhead
\bottomrule\noalign{}
\endlastfoot
\citep{robertson2025syn57} Syn57 & 4 Ser + 2 Ala + Stop & 60,240 &
Ser=100\%, Ala=0\% & Topology-fixed exposure (gene-level null) \\
\citep{fredens2019} Syn61 & Ser (TCG/TCA\(\rightarrow\)AGC/AGT) & 18,218
& All cross & Informative null \\
\citep{ostrov2016} & 7 codon types & 62,214 & Mixed & Segment
viability \\
\citep{napolitano2016} & Arg (AGR\(\rightarrow\)CGN) & 12,888 & None &
SRZ covariates \\
\citep{grome2025ochre} Ochre & Stop (TGA\(\rightarrow\)TAA) & 1,195 &
N/A & Growth data \\
\citep{nyerges2024}\textsuperscript{†} & 7 codon types & 62,007 & Mixed
& Multi-omics \\
\citep{lajoie2013} C321.\(\Delta\)A & Stop (UAG\(\rightarrow\)UAA) & 321
& N/A & Baseline \\
\citep{frumkin2018} & Arg (CGU/CGC\(\rightarrow\)CGG) & 60 & None &
Translation efficiency \\
\citep{ding2024} & Ser (TCG, mammalian) & Variable & Yes &
Cross-organism \\
\end{longtable}

\protect\phantomsection\label{tbl:codonsafe-datasets}{}

\begin{figure}
\centering
\emfig{Fig4\_translational}
\caption{Translational applications and prediction catalogue.
\textbf{(A)} Structural-preservation index (Supplement §S12) for 1,280
single-codon reassignments from the standard code, colored by
two-fold/four-fold filtration status. The index is a hand-weighted
discrete composite of four indicators and takes only four values across
the candidate universe (\(S \in \left\{ 0.55,0.75,0.80,1.00 \right\}\));
it is a visualization heuristic, not an inferential test. \textbf{(B)}
Grantham mismatch score decomposition by nucleotide position: position 2
dominates (49.3\%), consistent with the biochemical hierarchy of
mutational impact. \textbf{(C)} Prediction catalogue: 15 evaluated
claims grouped by workstream (WS1--WS4 plus WS6; WS5 is a synthesis
workstream that generates the catalogue itself and has no independent
claims), colored by claim-hierarchy category (Verified \(=\) Supported;
Tested \(=\) hypothesis-tested including Falsified; Exploratory;
Tautological; Falsified for the KRAS--Fano decision gate).}
\end{figure}

\protect\phantomsection\label{fig:translational}{}

\subsubsection{A three-layer interpretation of recoding outcomes
(exploratory)}\label{sec:res-threelayer}

This subsection is exploratory and hypothesis-generating. We propose a
three-layer view of how codon-space structure relates to recoding
outcomes, but the partition into evolutionary, recoding-burden, and
design-deviation layers was \textbf{not} prespecified before the
analyses below. The layering is therefore a working synthesis, not a
tested model: each contrast below is individually informative, but their
integration into a single three-layer schema is interpretive and would
require an independent dataset built around a sharper hypothesis to
confirm.

\paragraph{Boundary crossing: validated but not predictive of
transcriptomic perturbation}

In the \citep{robertson2025syn57} Syn57 dataset (the only large-scale
dataset with within-study variation in boundary crossing), all 37,146
serine recodings crossed the UCN\(\leftrightarrow\)AGY disconnection
boundary at \(\varepsilon = 1\), while all 22,859 alanine recodings
remained within the compact GCN family. This provides an exact
within-study topology contrast. However, genes containing exclusively
serine-type versus exclusively alanine-type recodings did not differ in
transcriptomic perturbation (Mann--Whitney \(p = 0.40\),
\(n_{\text{Ser }} = 129\), \(n_{\text{Ala }} = 49\)), and the fraction
of boundary-crossing recodings per gene was uncorrelated with
\(|\log_{2}\text{ FC}|\) (\(\rho = 0.012\), \(p = 0.46\), \(n =\)3,510
genes). In a separate dataset, all 18,218 Syn61 serine recodings
\citep{fredens2019} crossed the boundary; the seven design-to-final
corrections that the Syn61 authors required for cell viability
(Supplementary Data 19 of \citep{fredens2019}) concentrated at specific
functional units (a Ser70 site in \emph{yaaY}, an intergenic separation
between \emph{ftsI} and \emph{murE}, a Ser4 site in \emph{map}, five
positions in the \emph{yceQ} pseudogene needed to preserve viable
expression of the downstream essential gene \emph{rne}) rather than at
randomly distributed sites. Because all 18,218 forward swaps share the
same two AA-identity-preserving swap subtypes (TCG\(\rightarrow\)AGC and
TCA\(\rightarrow\)AGT), \textbar Δphys\textbar{} at the AA-identity
level is structurally constrained across all sites, and the 7
corrections do not lie at extreme \textbar Δphys\textbar{} or unusual
Hamming-accessibility positions relative to the 18,211 non-correction
sites. We therefore interpret the corrections as protein-level
functional pressure at specific residues and operon-context constraints,
not a global topology-driven burden, and treat this contrast as
descriptive rather than as a hypothesis-test of GF(2)\textsuperscript{6}
features. Together with the gene-level Syn57 null, this indicates that
boundary crossing is a coarse geometric property of codon-family
organization: validated as a structural distinction over evolutionary
timescales but not predictive of acute transcriptomic perturbation in
synthetic-biology engineering, where competing ecological pressures are
absent and expression constructs are highly optimized.

\paragraph{Local mismatch geometry: a fine-scale burden axis with
positive signal}\label{sec:local-mismatch-geom}

The second layer examines whether reassignment changes the
physicochemical quality of the immediate mutational neighborhood. In the
Syn57 contrast, serine swaps moved codons into better local Grantham
neighborhoods (\(\Delta F_{\text{local }} = - 37.2 \pm 30.5\)), whereas
alanine swaps moved codons into worse neighborhoods at a fixed
\(\Delta F_{\text{local }} = + 19.0 \pm 0.0\) (zero variance). This
contrast is design-confounded and we report it descriptively rather than
inferentially: every Syn57 alanine swap goes to the same target codon,
so all alanine \(\Delta_{\text{local}}\) values are identical, and the
Mann--Whitney \(U = 0\), \(p < 10^{- 16}\) comparison reflects
design-level non-overlap rather than a within-group inferential test.
The methodologically clean inferential test comes from a topology-fixed
setting (arginine synonymous recodings \citep{napolitano2016}, where all
12,888 AGR\(\rightarrow\)CGN swaps remain within one connected component
at \(\varepsilon = 1\)), in which local mismatch change correlated with
established recoding-burden covariates: \(\Delta F_{\text{local}}\) vs
RBS deviation (\(\rho = - 0.33\), \(p < 10^{- 15}\)) and vs mRNA folding
deviation (\(\rho = - 0.12\), \(p < 10^{- 15}\)). Among the four CGN
targets, CGG has the lowest local mismatch cost (246 Grantham units)
while CGC has the highest (421), and target-specific RBS deviation spans
four orders of magnitude, suggesting that local mismatch and SRZ
covariates capture partially independent dimensions of recoding burden.

\paragraph{Design deviations: accessibility-dominated, not
topology-directed}

To test whether engineered systems escape design constraints along
topology-defined low-burden directions, we compared the Syn57 design
genome (Data file S2; \citep{robertson2025syn57}) against the final
verified genome (Data file S8) and identified 727 CDS-level codon
differences. After filtering 7 genes with \(>\)10 deviations each
(structural rearrangements, predominantly \emph{bioF} with 311 changes),
326 genuine point deviations remained (267 synonymous, 59
nonsynonymous). Among the 163 deviations that changed local mismatch, 76
moved to a better neighborhood and 87 to a worse one, showing no
directional bias (two-sided exact binomial \(p = 0.43\),
Clopper--Pearson 95\% CI \(\lbrack 0.39,0.55\rbrack\); Wilcoxon
signed-rank \(p = 0.67\)). This null result was stable across all
filtering thresholds tested (\(p > 0.5\) at cutoffs 3, 5, 10, 15, 20,
and 50 deviations per gene). By contrast, deviations strongly favored
mutationally proximate moves: 83\% were single-bit changes (Hamming
distance 1 under the default encoding; note that Hamming distance is
encoding-dependent, though nucleotide edit distance is invariant), and
the mean Hamming distance was 1.21. Thus, realized deviations from the
Syn57 design do not preferentially descend toward lower local mismatch;
instead, they follow accessibility-dominated escape routes, favoring
nearby codon states regardless of neighborhood quality. The
\citep{ding2024} mammalian TCG recoding independently confirms the
universality of the serine disconnection across standard-code organisms,
providing cross-kingdom validation.

\paragraph{Segment-level recoding burden in the Ostrov 57-codon design}

As an orthogonal test, we examined segment-level outcomes from the
\citep{ostrov2016} 57-codon genome design. Among 44 segments with
fitness data (of 87 total), segments with more recoded codons in
essential genes showed a suggestive correlation with worse doubling time
(\(\rho = 0.34\), raw \(p = 0.022\), Bonferroni-corrected
\(p = 0.066\)), while total recoding load showed no association
(\(\rho = - 0.10\), \(p = 0.53\)). Problem segments (13 with lethal
exceptions) trended toward higher essential-gene recoding load (73.5 vs
39.8 recoded sites, Mann--Whitney \(p = 0.086\)). Among 44 fully
characterized segments, 17 showed spontaneous codon reversions,
indicating ongoing design instability at positions where the engineered
code was under selection pressure.

\subsubsection{Summary of cross-study reanalysis scope}

Across the four synthetic-recoding contrasts, the results are consistent
with (but do not prove) a three-layer view in which different structural
descriptors may operate at different biological layers (family topology,
local mismatch, and Hamming accessibility, respectively). Specific
testable predictions follow: (i) a Syn-style recoding experiment that
varies family-boundary crossing while holding local mismatch fixed
should fail to predict transcriptomic burden (already partly observed in
Syn57); (ii) a recoding scan that varies \(\Delta F_{\text{local}}\)
within a single boundary regime should predict RBS-context proxies
(consistent with the Napolitano arginine result); (iii) directed
evolution starting from a recoded design should produce escape mutations
whose Hamming distribution is shaped more by local accessibility than by
mismatch optimisation (consistent with the Syn57 design-deviation
analysis). Each prediction admits a direct experimental test that would
either confirm or falsify the layer structure proposed here.

\subsection{Scope of the framework: rejected and null
findings}\label{sec:res-negatives}

Four conjectural extensions of the \({\text{GF}(2)}^{6}\) representation
were tested and cleanly separated from the surviving picture: (i) the
KRAS--Fano clinical prediction was falsified against 1,670 MSK-IMPACT
\citep{zehir2017} mutations (\(p = 1.0\); Supplement §S13); (ii) the
claim that Serine's inter-family Hamming distance equals 4 under all 24
base-to-bit encodings is false (16 encodings yield distance 2 and only 8
yield distance 4), though Serine is disconnected at \(\varepsilon = 1\)
under every encoding (Supplement §S14); (iii) PSL(2,7) has no
64-dimensional irreducible linear representation and is therefore
pre-rejected as a symmetry group under that specific realisation
\citep{antoneli2011} (this narrower claim does not exclude other group
actions on a 64-element set; Supplement §S15); and (iv) the
coordinate-wise map \({\text{GF}(2)}^{6} \rightarrow {\mathbb{C}}^{3}\)
sending base pairs to fourth roots of unity fails the required
group-character identity \(\chi(x + x) = {\chi(x)}^{2}\) since
\(i^{2} = - 1 \neq 1\); the ``holomorphic'' qualifier adds no analytic
content on a finite discrete domain (Supplement §S16). A separate null
result on source-neighborhood burden (Mann--Whitney \(U = 301\),
\(p = 0.70\); Supplement §S17) indicates that reassignment is not driven
by local escape from costly source neighborhoods, and clarifies that the
topology-avoidance constraint (§3.4) operates at the global
graph-connectivity level rather than per-codon.

\section{Discussion}\label{sec:discussion}

\subsection{An information-theoretic view of the genetic
code}\label{sec:disc-info}

The central finding of this work is that the standard genetic code
minimizes the physicochemical disruption caused by single-bit errors in
\({\text{GF}(2)}^{6}\) coordinates. This is not a new conclusion
(\citep{freeland1998} established error-minimization using
nucleotide-level mutation models), but the hypercube representation
makes the optimality principle geometrically explicit: the code is a
good \emph{coloring} of a structured graph, in the graph-theoretic sense
that adjacent vertices (codons differing by one bit) tend to share
labels (amino acids) or, when they differ, differ by small
physicochemical distances.

The score decomposition (\Cref{fig:translational}, panel B)
concentrates optimization at the second (49.3\%) and first (38.2\%)
codon positions, with the wobble position contributing only 12.5\% ---
an emergent property of the code's structure rather than a model
parameter. The \(\rho\)-robustness result (\Cref{fig:coloring}, panel
B) demonstrates that optimality is not an artifact of restricting
attention to \(Q_{6}\): when the full \(H(3,4)\) mutation graph is
considered (\(\rho = 1\)), the signal strengthens. This
error-minimization is complementary to, but distinct from, the finding
of \citep{itzkovitz2007} that the code is also nearly optimal for
carrying parallel regulatory information within protein-coding
sequences.

The \({\text{GF}(2)}^{6}\) representation contributes two analytical
tools beyond a pure \(H(3,4)\) analysis. First, the coordinatized
decomposition of \(H(3,4)\) into the 192 Hamming-1 edges of \(Q_{6}\)
plus 96 within-nucleotide diagonal (Hamming-2) edges enables the
weighted-mismatch score \(F_{\rho}\) (\Cref{eq:weighted}) to
interpolate continuously from \(\rho = 0\) (\(Q_{6}\) only) to
\(\rho = 1\) (full \(H(3,4)\)); the monotonic strengthening of
optimality as \(\rho \rightarrow 1\) (Table 3, §3.2) does not exist in a
pure \(H(3,4)\) analysis, where the diagonal partition is not
coordinatized. Second, the Hamming-distance filtration
\(\beta_{0}\left( G_{a}^{\varepsilon} \right)\) provides a natural
persistence parameter for the disconnection catalogue
(\Cref{fig:topology}, panel A) and for the conditional-logit topology
feature (\(\Delta\beta_{0}\)). The 24-encoding sensitivity audit is a
necessary control rather than a contribution: it identified \(Q_{6}\)
topology depletion as encoding-dependent (8 of 24 bijections give no
signal; §3.4, Supplement §S4) and prompted promotion of \(H(3,4)\) to
primary, but does not extend the framework's reach. Pure \(H(3,4)\) has
no encoding choice and needs no such audit.

The binary representation is therefore best understood as an analytical
decomposition tool; its distinct contributions are the \(\rho\)-sweep
continuum and the \(\varepsilon\)-filtration, with \(\rho\) a
diagonal-edge weight rather than a transition/transversion weight
(Methods §2.1, Supplement §S2).

\subsection{Evolutionary preservation and topology
avoidance}\label{sec:disc-evol}

The per-table analysis (\Cref{fig:coloring}) shows that 26 of 27 NCBI
translation tables maintain coloring optimality under their own
quartet-pattern shuffle null, suggesting that codon reassignment events
are constrained to preserve error-minimization. The marginal exception
is the most extensively reassigned code: translation table 3 (yeast
mitochondrial, 6 codon changes) falls only slightly above the 5\%
threshold.

The topology avoidance result (\Cref{fig:topo}) provides a mechanistic
complement: natural reassignments avoid creating new amino acid
disconnections at a rate far below chance expectation. This depletion is
not confined to the \(Q_{6}\) representation: under the full \(H(3,4)\)
single-nucleotide mutation graph, the observed proportion of
topology-breaking reassignments remains 21.4\% against a possible rate
of 66.1\% (\(\text{RR } = 0.32\), 95\% CI \(\lbrack 0.16,0.66\rbrack\),
\(p \leq 10^{- 4}\)). The signal magnitude is similar to the \(Q_{6}\)
result (\(\text{RR } = 0.29\)); some \(Q_{6}\)-disconnected pairs become
connected when all single-nucleotide edges are admitted, slightly
lowering the possible-rate denominator, but the qualitative conclusion
survives: codon-family connectivity is functionally important under
either the hypercube subgraph or the fuller single-substitution graph.
The yeast mitochondrial threonine reassignment illustrates this cost:
the CUN\(\rightarrow\)Thr change required acquisition of a novel
tRNA\textsuperscript{Thr} derived from tRNA\textsuperscript{His} via
anticodon mutation \citep{su2011}, creating the topology-breaking
disconnection that makes translation table 3 the sole marginal outlier
in the per-table optimality analysis.

Together these analyses indicate that code evolution is constrained
along two partly independent axes, not one. Physicochemical smoothness
preserves \textbf{protein function under mistranslation}: a reassignment
that lands in a physicochemically similar neighborhood is one whose
mistranslation errors remain tolerable. Topological integrity preserves
\textbf{something different}: the connectivity of an amino-acid codon
family determines whether existing decoding machinery (a tRNA serving
wobble-related codons) can continue to service the family without new
molecular hardware. The conditional logit analysis
(\Cref{tbl:condlogit}) makes the separation explicit. The topology term
retains major explanatory value (\(\Delta\text{AICc } = 108\) relative
to a physicochemistry-only model) after accounting for local
physicochemical cost, while physicochemical cost retains comparable
value (\(\Delta\text{AICc } = 89\)) after accounting for topology; the
two feature classes are only weakly correlated across the candidate
landscape (\(r_{s} = 0.15\)), so neither term is a redundant restatement
of the other.

To rule out a coordinate artifact, the M3 refit with the
encoding-independent \(H(3,4)\) topology feature retains a similar
explanatory value
(\(\Delta\text{AICc}\left( M1\  \rightarrow \ M3{H(3,4)} \right) = 91\)
vs. \(\Delta\text{AICc}(M1\  \rightarrow \ M3Q\_ 6) = 108\); §3.5),
confirming that the conditional-logit thesis does not depend on the
choice of topology graph.

We do not interpret this as direct selection on topology as an abstract
graph property; topology may still proxy aspects of decoding
architecture (tRNA cross-recognition range, ribosomal A-site
constraints) not explicitly modeled here. The central claim is
structural: within the tested event landscape, topology avoidance cannot
be dismissed as a by-product of physicochemical optimization or as a
coordinate artifact of the binary representation. The contribution
beyond single-axis treatments is the second axis itself: code evolution
moves through a low-cost region \textbf{and} avoids fragmentation of the
decoding substrate, with the two constraints partly orthogonal and
neither reducible to the other.

\subsection{Exploratory tRNA accommodation across variant-code
systems}\label{sec:disc-trna}

The tRNA enrichment result (\Cref{fig:topo}, panel D;
\Cref{tbl:trna-tests}) describes an \textbf{accommodation-of-decoding}
pattern across variant-code lineages: variant-code genomes tend to carry
elevated tRNA-gene counts for the reassigned amino acid relative to
standard-code sisters (median MIS Stouffer \(p = 0.037\) across 332 MIS
of size 6). The signal is driven principally by
topology-\textbf{preserving} stop-to-sense reassignments ---
UAR\(\rightarrow\)Gln in ciliates, UGA\(\rightarrow\)Cys in
\emph{Euplotes}, UGA\(\rightarrow\)Trp in \emph{Blepharisma} and
\emph{Mycoplasmoides} --- rather than by the four
topology-\textbf{breaking} pairings, whose subset is null (Stouffer
\(p = 0.387\)). The result does not therefore establish compensation for
codon-family graph disconnection; it is consistent with a broader
``accommodation of decoding'' whose distribution across lineages is
heterogeneous. Case examples illustrate the variety of mechanisms in
this heterogeneous set: \emph{Tetrahymena thermophila} carries 54 Gln
tRNAs including a large suppressor complement; \emph{Blastocrithidia
nonstop} achieves UGA\(\rightarrow\)Trp via anticodon stem shortening
rather than tRNA-gene duplication \citep{kachale2023};
\emph{Mycoplasmoides} species with UGA\(\rightarrow\)Trp read both
codons with a single tRNA-Trp through anticodon modification. The
correlation between the mechanistic route (duplication vs anticodon
modification vs base modification) and genome size that these three
cases motivate is a \textbf{hypothesis} about how genome size might
constrain the mechanistic repertoire, not a statistical association test
in this manuscript.

\subsection{Synthetic recoding outcomes: a three-layer
interpretation}\label{sec:disc-layers}

The retrospective cross-study reanalysis helps delimit which structural
quantities in GF(2)\textsuperscript{6} are functionally relevant and at
which biological layer. Boundary crossing is a valid codon-space
distinction (serine and alanine provide a clean contrast in which one
class always crosses a family boundary and the other never does), yet
this distinction did not predict RNA-seq perturbation in Syn57. We
therefore do not interpret codon-family boundary crossing as a general
transcriptomic burden variable. Its relevance, if any, likely lies at a
different biological layer, such as decoding architecture, translational
fidelity, or evolutionary accessibility, rather than steady-state
expression change.

By contrast, local mismatch geometry showed a clearer positive signal.
In both the Syn57 contrast and the topology-fixed arginine setting,
codon changes differed systematically in the quality of the
physicochemical neighborhood they entered, and these differences tracked
established recoding-burden covariates \citep{napolitano2016}. This
suggests that local mismatch is best understood as a fine-scale
recoding-friction variable: not a universal master predictor, but a
meaningful descriptor of neighborhood quality within an already
specified codon-family structure.

The design-deviation analysis further clarifies scope. Genuine
deviations from the Syn57 design were overwhelmingly short-range (83\%
single-bit in GF(2)\textsuperscript{6}), with no directional bias in
mismatch change. This indicates that realized escape is governed more by
mutational accessibility than by systematic optimization of local
neighborhood cost: once a design is under strain, the available exits
appear to be chosen primarily by mutational proximity.

These observations are consistent with, but do not establish, a working
hypothesis that ``topology'' should not be treated as a single
explanatory variable. Under that hypothesis, different structural
descriptors would matter at different biological layers: family topology
for the geometry of reassignment space (\Cref{sec:res-topo}), local
mismatch for some aspects of recoding burden
(\Cref{sec:local-mismatch-geom}), and Hamming accessibility for the
short-range routes by which engineered designs deviate. Each of the
three layer claims rests on a single retrospective contrast and the
partition into layers was not prespecified, so the schema should be read
as hypothesis-generating rather than as a tested model (see §3.8.2 for
explicit predictions). That the \citep{fredens2019} Syn61 design
tolerated 18,218 boundary-crossing serine swaps as a class (a
genome-wide engineering success rate, not a per-codon-position viability
test), while the NCBI variant codes show 3.1-fold depletion of
topology-breaking changes over evolutionary timescales, is consistent
with this layered reading: boundary crossing may be more relevant to
evolutionary trajectory constraints than to acute engineering costs, but
a within-study experiment that varies family-boundary crossing while
holding local mismatch fixed is required to test that prediction.

\subsection{Broader connections}\label{sec:disc-broader}

The \({\text{GF}(2)}^{6}\) framework connects to three adjacent
literatures that we treat only briefly here. A complementary line of
BioSystems work treats the code algebraically using 2-adic dynamical
systems on \(\mathbb{Z}_{2}\)
\citep{khrennikov2007, dragovich2019, axelsson2024a, axelsson2024b}. The
two descriptions use different geometries (ultrametric attractors vs.
finite Hamming graphs) and we do not identify them. As a finite bridge,
the Walsh--Hadamard / 2-adic spectral depth of the standard code (544,
encoding-invariant across all 24 bijections) sits far below a block-size
matched null mean of \(689.3 \pm 8.2\) at \(n =\)2,000
(\(z = - 17.74\)), but is mathematically constant under a stricter
wobble-box-preserving label-permutation null; the Walsh signature is
therefore an algebraic invariant of the wobble-box structure rather than
a separate beyond-wobble optimality signal (Supplement §S21; the
complementary wobble-free label-spectrum fraction \(S = 0.7514\) is
included as a second signature of the same structure).

A second line of work by \citep{gonzalez2016} proposes a
\textbf{non-power} binary representation of the code based on a
redundant signed-binary numeration with Fibonacci-like positional
weights \(\lbrack 1,1,2,4,7,8\rbrack\), in which codon-to-integer maps
become intentionally many-to-one, mirroring code degeneracy. The 24
base-to-bit bijections we sweep are all \textbf{power} representations,
in the sense that each codon's 6-bit vector is the concatenation of
three per-base 2-bit encodings; the non-power family is strictly larger
and would give a different hypercube-like object with an edge structure
not factorable across the three codon positions. Our primary
topology-avoidance result is defined on the encoding-independent
\(H(3,4)\) mutation graph (nucleotide-level, invariant under any
codon-level bijection whatsoever), so the paper's principal claim is not
exposed to this concern. The \(Q_{6}\) subgraph analysis, which is
already encoding-dependent within the power family (§3.4, §S4), would
benefit from a non-power extension; we flag this as a natural
continuation but do not attempt the full sweep here.

Separately, \citep{tsour2026}`s deep proteomic survey of alternate RNA
decoding in over 1,000 human samples identified 8,746 unique amino-acid
substitutions; among the 5,611 unambiguous-pair events, 65.0\% involve
substitutions whose closest source--target codon pair differs at a
single nucleotide, versus a baseline of 39.5\%. The naive event-level
binomial statistic ignores clustering of substitutions by substitution
type, gene, sample, and measurement opportunity; the effect-size
contrast (65.0\% vs 39.5\%) is the primary summary, and type-level
analyses are given in Supplement §S22. This provides external empirical
grounding for the codon-adjacency mechanism implicit in our
edge-mismatch objective (\Cref{eq:mismatch}): the same
single-nucleotide neighbourhood structure governs the observed
alternate-translation events in mammalian proteomes, independent of the
amino-acid distance metric.

\subsection{Limitations}\label{sec:disc-limits}

Several limitations should be noted.

\textbf{Encoding choice.} The base-to-bit encoding is not unique. While
coloring optimality holds across all 24 encodings, specific score values
and rank orderings are encoding-dependent (Supplementary Material).

\textbf{tRNA enrichment fragility.} The tRNA enrichment signal is
present in the median-independent-subset sense (median Stouffer
\(p = 0.037\) across 332 MIS of size 6) but does not survive the strict
worst-case subset criterion (worst-case \(p = 0.104\)), and the
pre-specified topology-breaking subset (\(n = 4\)) is null (Stouffer
\(p = 0.387\)). We therefore classify the result as exploratory rather
than as demonstrating a universal or mechanism-specific compensation for
codon-family disconnection. Assembly-fragmentation robustness (removing
two genomes with \(> 10\)\% pseudogene fraction yields Stouffer
\(p = 0.041\) on the remaining 22 pairings) is reported in Supplement
§S10.

\textbf{Conditional logit scope.} The conditional logit model
(\Cref{sec:res-condlogit}) is event-level explanatory, not direct proof
of biological causality. The candidate universe comprises all
\(\approx\)1,280 single-codon reassignments regardless of mechanistic
plausibility, and the tRNA-complexity proxy (Hamming distance to nearest
target-AA codon) is heuristic. The non-significance of this particular
proxy does not exclude that a richer model of tRNA repertoire change
would contribute explanatory value.

\textbf{Survivorship bias.} All event-level analyses operate on
reassignments that \emph{persisted} in extant lineages.
Topology-breaking reassignments that produced fitness collapse and
lineage extinction are unobservable. The depletion result is therefore
consistent with both selection against attempting topology-breaking
moves and selection against the lineages that attempted them;
cross-sectional NCBI data cannot adjudicate. Both interpretations
support the conclusion that codon-family connectivity constrains
evolutionary trajectories, with different mechanistic implications for
the locus of selection.

\textbf{Phylogenetic non-independence ceiling.} The 28 de-duplicated
reassignment events used in the topology-avoidance hypergeometric test
cluster at the codon level (e.g., the recurrent UGA\(\rightarrow\)Trp
event is observed across many distantly related mitochondrial lineages
and is collapsed to a single event for the de-duplicated test) but they
are \textbf{not} independent at the level of evolutionary origin:
mitochondrial codes share an ancestral mitochondrial-decoding regime,
ciliate nuclear codes share a ciliate-specific eRF1 trajectory, and so
on. A conservative lower bound on the number of phylogenetically
independent topology-preservation origins is on the order of 4--6, well
below the 22 non-breakers the de-duplicated count reports; the
hypergeometric \(p\) should be read against that ceiling rather than
against \(n = 22\). The conditional-logit clade-exclusion sensitivity
(Supplement §S7) and the within-clade enumerations in §3.4 jointly bound
how strongly the depletion magnitude can be claimed.

Multiple-comparison correction was applied within analysis families
rather than across all quantities; the eight test families address
conceptually distinct, non-nested questions, and all headline results
survive cross-family Bonferroni at \(\alpha = \frac{0.05}{8}\) except
the exploratory tRNA enrichment (detailed breakdown in Supplement §S18).

\section{Conclusion}\label{sec:conclusion}

The standard genetic code is consistently low-cost: across four amino
acid distance metrics and across a continuum of mutation graphs spanning
\(Q_{6}\) and the complete single-nucleotide graph \(H(3,4)\), it sits
in the extreme tail of quartet-pattern shuffle null models
(\(p \leq 0.006\) under every combination). Across the 27 NCBI
translation tables this near-optimality is preserved at the variant
level: of the 12 codes whose distance from the standard makes the
per-table test informative, 11 retain top-5\% placement after BH--FDR
correction, with only yeast mitochondrial (table 3) marginally above the
5\% threshold; the 14 near-standard tables are formally significant in
the same correction but we interpret these as primarily confirmatory of
standard-code geometry. Code evolution therefore appears to occur within
a constrained low-cost region rather than freely across code space.

Two partly independent constraints govern code evolution within that
region. Persistent natural reassignments preferentially enter
neighborhoods of lower local Grantham mismatch, and they avoid
fragmenting amino-acid codon-family connectivity. Topology-breaking
moves are approximately 3.1-fold depleted under the encoding-independent
\(H(3,4)\) adjacency, and conditional-logit decomposition shows that
topology adds explanatory value beyond physicochemistry, and vice versa
(\(\Delta\text{AICc } \geq 89\) under \(Q_{6}\) topology;
\(\Delta\text{AICc } \geq 91\) under encoding-independent \(H(3,4)\)
topology). The two axes are only weakly correlated across the candidate
landscape (\(r_{s} = 0.15\)); neither is a redundant restatement of the
other. The topology coefficient survives all seven phylogenetic clade
exclusions (\(\Delta\text{AICc}(M1\  \rightarrow \ M3) \geq 49\) in the
worst regime), and a parametric predictive simulation under M3
reproduces the observed event-step topology-breaking rate under
\(Q_{6}\) / \(\Delta\beta_{0} > 0\) (5/66 = 0.076 vs simulated mean
0.077; \(p = 0.60\)). Several variant-code lineages show exploratory
tRNA-gene enrichment for the reassigned amino acid (median MIS Stouffer
\(p = 0.037\) across 24 pairings among 24 organisms of mixed provenance;
worst-case \(p = 0.104\); topology-breaking subset \(p = 0.39\)),
consistent with heterogeneous accommodation-of-decoding rather than a
specific compensation-for-disconnection mechanism.

A retrospective cross-study reanalysis of eight published
genome-recoding experiments (\(>\)217,000 codon-level observations)
sharpens the scope of the topology constraint. Family-boundary crossing
is a real structural distinction over evolutionary timescales but did
\textbf{not} predict acute transcriptomic perturbation in the Syn57
gene-level contrast, while Syn61 had earlier tolerated 18,218
boundary-crossing serine swaps as a class \citep{fredens2019} (a
genome-wide success rate, not a per-codon-position viability measure).
This pattern is consistent with topology acting primarily as a
constraint on evolutionary trajectories: gradual fragmentation of a
codon family could create periods of decoding fragility that are
bypassed in synthetic systems where decoding machinery is engineered
up-front. Codon-space topology may therefore be a constraint on how
genetic codes can change rather than on how they currently function:
path-dependent rather than instantaneous. The three-layer integration of
these descriptors (evolutionary topology, fine-scale recoding burden,
accessibility-driven escape) is exploratory and hypothesis-generating:
each layer-claim rests on a single retrospective contrast, the partition
was not pre-specified, and confirmation requires the prospective tests
outlined in §3.8.2.

The \({\text{GF}(2)}^{6}\) representation is best understood as an
analytical decomposition rather than as a claim about biologically
privileged binary coordinates. Its main contributions are the
\(\rho\)-sweep interpolating between \(Q_{6}\) and \(H(3,4)\), the
\(\varepsilon\)-filtration of codon-family connectivity, and the
encoding-sensitivity audit required for representation-dependent
analyses. The KRAS--Fano clinical prediction, Serine distance-4
invariant, PSL(2,7) symmetry, and holomorphic-embedding claims were
tested and rejected (\Cref{tbl:claims}): several invariants of the
representation turned out to be encoding artifacts. The surviving
picture is encoding-independent: the genetic code is a low-cost coloring
of \(H(3,4)\), and observed reassignment trajectories preserve both
physicochemical similarity and codon-family connectivity. Whether the
topology axis reflects direct selection or a proxy for a specific
feature of decoding architecture remains open. The present results
support a working model in which natural genetic-code evolution is
constrained by two partly independent axes: physicochemical smoothness
and codon-family connectivity.

\section{Acknowledgements}

We thank the NCBI, GtRNAdb, and cBioPortal teams for maintaining public
databases essential to this work. tRNAscan-SE 2.0.12 was developed by
Chan and Lowe at UC Santa Cruz.

\section{Funding}

This research did not receive any specific grant from funding agencies
in the public, commercial, or not-for-profit sectors.

\section{CRediT author contribution statement}

\textbf{Paul Clayworth:} Conceptualization, Methodology, Formal
analysis, Writing -- original draft, Writing -- review and editing.
\textbf{Sergey Kornilov:} Conceptualization, Methodology, Software,
Formal analysis, Data curation, Visualization, Validation, Writing --
original draft, Writing -- review and editing.

\section{Declaration of generative AI and AI-assisted technologies}

During the preparation of this work the authors used GPT-5.2-Pro,
GPT-5.5-Pro, Claude Opus 4.5, Claude Opus 4.7, GLM-5.1, Kimi-K2.5,
MiniMax-M2.7, and Gemini 3 Pro to scaffold and review code, review the
manuscript outline, and assist in drafting and revising the manuscript
text. After using these tools, the authors reviewed and edited the
content as needed and take full responsibility for the content of the
publication. No generative AI or AI-assisted tools were used to create
or alter figures or images in this manuscript.

\section{Declaration of competing interest}

The authors declare no competing interests.

\section{Ethical statement}

This study did not involve human subjects, animal experiments, or
clinical samples. All analyses were performed on publicly available
data: NCBI genome assemblies, NCBI gc.prt translation table definitions,
GtRNAdb tRNA gene catalogues, the MSK-IMPACT mutation registry as
released through cBioPortal, and previously published genome-recoding
datasets. No new biological materials were generated.

\section{Data and code availability}

All code, raw data, and intermediate analysis outputs are publicly
released in the \texttt{codontopo} repository at
\url{https://github.com/biostochastics/codontopo} (version 0.6.1, tag
\texttt{v0.6.1}). Install with \texttt{pip\ install\ -e\ ".{[}all{]}"}
or \texttt{uv\ sync\ -\/-all-extras}. Analyses are fully reproducible
via:

\begin{verbatim}
git clone https://github.com/biostochastics/codontopo.git
cd codontopo
codon-topo all --output-dir=./output --seed=135325
\end{verbatim}

All inline statistics, tables, and figures in this manuscript and
supplement are rendered directly from the pipeline outputs by the Typst
sources (\texttt{manuscript.typ}, \texttt{supplement.typ}) included in
the same repository. NCBI genome assembly accessions are listed in
Supplement §S10.


% Corrected 2026-08-15: bibliography inlined directly (rather than
% relying on \bibliography{references} + a separately-shipped
% manuscript.bbl) so the reference list is embedded in this single
% .tex file regardless of how Editorial Manager tags/stages
% auxiliary files. Content below is the bibtex-generated .bbl
% (elsarticle-harv style) verbatim.
\begin{thebibliography}{40}
\expandafter\ifx\csname natexlab\endcsname\relax\def\natexlab#1{#1}\fi
\providecommand{\url}[1]{\texttt{#1}}
\providecommand{\href}[2]{#2}
\providecommand{\path}[1]{#1}
\providecommand{\DOIprefix}{doi:}
\providecommand{\ArXivprefix}{arXiv:}
\providecommand{\URLprefix}{URL: }
\providecommand{\Pubmedprefix}{pmid:}
\providecommand{\doi}[1]{\href{http://dx.doi.org/#1}{\path{#1}}}
\providecommand{\Pubmed}[1]{\href{pmid:#1}{\path{#1}}}
\providecommand{\bibinfo}[2]{#2}
\ifx\xfnm\relax \def\xfnm[#1]{\unskip,\space#1}\fi
%Type = Article
\bibitem[{Antoneli and Forger(2011)}]{antoneli2011}
\bibinfo{author}{Antoneli, F.}, \bibinfo{author}{Forger, M.},
  \bibinfo{year}{2011}.
\newblock \bibinfo{title}{Symmetry breaking in the genetic code: finite
  groups}.
\newblock \bibinfo{journal}{Mathematical and Computer Modelling}
  \bibinfo{volume}{53}, \bibinfo{pages}{1469--1488}.
\newblock \DOIprefix\doi{10.1016/j.mcm.2010.03.050}.
%Type = Article
\bibitem[{Benjamini and Hochberg(1995)}]{benjamini1995}
\bibinfo{author}{Benjamini, Y.}, \bibinfo{author}{Hochberg, Y.},
  \bibinfo{year}{1995}.
\newblock \bibinfo{title}{Controlling the false discovery rate: a practical and
  powerful approach to multiple testing}.
\newblock \bibinfo{journal}{Journal of the Royal Statistical Society: Series B
  (Methodological)} \bibinfo{volume}{57}, \bibinfo{pages}{289--300}.
\newblock \DOIprefix\doi{10.1111/j.2517-6161.1995.tb02031.x}.
%Type = Book
\bibitem[{Burnham and Anderson(2002)}]{burnham2002}
\bibinfo{author}{Burnham, K.P.}, \bibinfo{author}{Anderson, D.R.},
  \bibinfo{year}{2002}.
\newblock \bibinfo{title}{Model Selection and Multimodel Inference: A Practical
  Information-Theoretic Approach}.
\newblock \bibinfo{edition}{2nd} ed., \bibinfo{publisher}{Springer}.
\newblock \DOIprefix\doi{10.1007/b97636}. \bibinfo{note}{source of the
  conventional $\Delta$AIC > 10 evidence threshold; calibrated on
  linear-regression model-comparison contexts.}
%Type = Article
\bibitem[{Buschmann et~al.(2026)Buschmann, Bolz, El-Hendi, Malekian and
  Schroeder}]{buschmann2026}
\bibinfo{author}{Buschmann, L.}, \bibinfo{author}{Bolz, S.N.},
  \bibinfo{author}{El-Hendi, F.}, \bibinfo{author}{Malekian, N.},
  \bibinfo{author}{Schroeder, M.}, \bibinfo{year}{2026}.
\newblock \bibinfo{title}{Much ado about nothing: modeling amino acid
  replacement with predicted protein structures}.
\newblock \bibinfo{journal}{Bioinformatics} \bibinfo{volume}{42},
  \bibinfo{pages}{btag188}.
\newblock \DOIprefix\doi{10.1093/bioinformatics/btag188}.
%Type = Article
\bibitem[{Chan and Lowe(2019)}]{chan2019}
\bibinfo{author}{Chan, P.P.}, \bibinfo{author}{Lowe, T.M.},
  \bibinfo{year}{2019}.
\newblock \bibinfo{title}{{tRNAscan-SE}: searching for {tRNA} genes in genomic
  sequences}.
\newblock \bibinfo{journal}{Methods in Molecular Biology}
  \bibinfo{volume}{1962}, \bibinfo{pages}{1--14}.
\newblock \DOIprefix\doi{10.1007/978-1-4939-9173-0_1}.
%Type = Unpublished
\bibitem[{Clayworth(2026)}]{clayworth2026}
\bibinfo{author}{Clayworth, P.}, \bibinfo{year}{2026}.
\newblock \bibinfo{title}{Algebraic structure of the genetic code: a technical
  note}.
\newblock \bibinfo{note}{Companion methodological note. Source code and full
  claim hierarchy are released alongside the present manuscript at
  https://github.com/biostochastics/codontopo; the published manuscript
  supersedes the technical note in scope and statistical detail.}
%Type = Article
\bibitem[{Cock et~al.(2009)Cock, Antao, Chang, Chapman, Cox, Dalke, Friedberg,
  Hamelryck, Kauff, Wilczynski and de~Hoon}]{cock2009}
\bibinfo{author}{Cock, P.J.A.}, \bibinfo{author}{Antao, T.},
  \bibinfo{author}{Chang, J.T.}, \bibinfo{author}{Chapman, B.A.},
  \bibinfo{author}{Cox, C.J.}, \bibinfo{author}{Dalke, A.},
  \bibinfo{author}{Friedberg, I.}, \bibinfo{author}{Hamelryck, T.},
  \bibinfo{author}{Kauff, F.}, \bibinfo{author}{Wilczynski, B.},
  \bibinfo{author}{de~Hoon, M.J.L.}, \bibinfo{year}{2009}.
\newblock \bibinfo{title}{{Biopython}: freely available {Python} tools for
  computational molecular biology and bioinformatics}.
\newblock \bibinfo{journal}{Bioinformatics} \bibinfo{volume}{25},
  \bibinfo{pages}{1422--1423}.
\newblock \DOIprefix\doi{10.1093/bioinformatics/btp163}.
%Type = Article
\bibitem[{Di~Giulio(2001)}]{digiulio2001}
\bibinfo{author}{Di~Giulio, M.}, \bibinfo{year}{2001}.
\newblock \bibinfo{title}{The origin of the genetic code cannot be studied
  using measurements based on the {PAM} matrix because this matrix reflects the
  code itself, making any such analyses tautologous}.
\newblock \bibinfo{journal}{Journal of Theoretical Biology}
  \bibinfo{volume}{208}, \bibinfo{pages}{141--144}.
\newblock \DOIprefix\doi{10.1006/jtbi.2000.2206}.
%Type = Article
\bibitem[{Ding et~al.(2024)Ding, Yu, Chen, Lao, Fang, Fang, Zhao, Yang and
  Lin}]{ding2024}
\bibinfo{author}{Ding, W.}, \bibinfo{author}{Yu, W.}, \bibinfo{author}{Chen,
  Y.}, \bibinfo{author}{Lao, L.}, \bibinfo{author}{Fang, Y.},
  \bibinfo{author}{Fang, C.}, \bibinfo{author}{Zhao, H.},
  \bibinfo{author}{Yang, B.}, \bibinfo{author}{Lin, S.}, \bibinfo{year}{2024}.
\newblock \bibinfo{title}{Rare codon recoding for efficient noncanonical amino
  acid incorporation in mammalian cells}.
\newblock \bibinfo{journal}{Science} \bibinfo{volume}{384},
  \bibinfo{pages}{1134--1142}.
\newblock \DOIprefix\doi{10.1126/science.adm8143}.
%Type = Article
\bibitem[{Dragovich and Mi{\v s}i{\'c}(2019)}]{dragovich2019}
\bibinfo{author}{Dragovich, B.}, \bibinfo{author}{Mi{\v s}i{\'c}, N.{\v Z}.},
  \bibinfo{year}{2019}.
\newblock \bibinfo{title}{p-{A}dic hierarchical properties of the genetic
  code}.
\newblock \bibinfo{journal}{BioSystems} \bibinfo{volume}{185},
  \bibinfo{pages}{104017}.
\newblock \DOIprefix\doi{10.1016/j.biosystems.2019.104017}.
%Type = Article
\bibitem[{Fredens et~al.(2019)Fredens, Wang, de~la Torre, Funke, Robertson,
  Christova, Chia, Schmied, Dunkelmann, Ber{\'a}nek, Uttamapinant, Llamazares,
  Elliott and Chin}]{fredens2019}
\bibinfo{author}{Fredens, J.}, \bibinfo{author}{Wang, K.},
  \bibinfo{author}{de~la Torre, D.}, \bibinfo{author}{Funke, L.F.H.},
  \bibinfo{author}{Robertson, W.E.}, \bibinfo{author}{Christova, Y.},
  \bibinfo{author}{Chia, T.}, \bibinfo{author}{Schmied, W.H.},
  \bibinfo{author}{Dunkelmann, D.L.}, \bibinfo{author}{Ber{\'a}nek, V.},
  \bibinfo{author}{Uttamapinant, C.}, \bibinfo{author}{Llamazares, A.G.},
  \bibinfo{author}{Elliott, T.S.}, \bibinfo{author}{Chin, J.W.},
  \bibinfo{year}{2019}.
\newblock \bibinfo{title}{Total synthesis of {Escherichia coli} with a recoded
  genome}.
\newblock \bibinfo{journal}{Nature} \bibinfo{volume}{569},
  \bibinfo{pages}{514--518}.
\newblock \DOIprefix\doi{10.1038/s41586-019-1192-5}.
%Type = Article
\bibitem[{Freeland and Hurst(1998)}]{freeland1998}
\bibinfo{author}{Freeland, S.J.}, \bibinfo{author}{Hurst, L.D.},
  \bibinfo{year}{1998}.
\newblock \bibinfo{title}{The genetic code is one in a million}.
\newblock \bibinfo{journal}{Journal of Molecular Evolution}
  \bibinfo{volume}{47}, \bibinfo{pages}{238--248}.
\newblock \DOIprefix\doi{10.1007/PL00006381}.
%Type = Article
\bibitem[{Frumkin et~al.(2018)Frumkin, Lajoie, Gregg, Hornung, Church and
  Pilpel}]{frumkin2018}
\bibinfo{author}{Frumkin, I.}, \bibinfo{author}{Lajoie, M.J.},
  \bibinfo{author}{Gregg, C.J.}, \bibinfo{author}{Hornung, G.},
  \bibinfo{author}{Church, G.M.}, \bibinfo{author}{Pilpel, Y.},
  \bibinfo{year}{2018}.
\newblock \bibinfo{title}{Codon usage of highly expressed genes affects
  proteome-wide translation efficiency}.
\newblock \bibinfo{journal}{Proceedings of the National Academy of Sciences}
  \bibinfo{volume}{115}, \bibinfo{pages}{E4940--E4949}.
\newblock \DOIprefix\doi{10.1073/pnas.1719375115}.
%Type = Article
\bibitem[{Gonzalez et~al.(2016)Gonzalez, Giannerini and Rosa}]{gonzalez2016}
\bibinfo{author}{Gonzalez, D.L.}, \bibinfo{author}{Giannerini, S.},
  \bibinfo{author}{Rosa, R.}, \bibinfo{year}{2016}.
\newblock \bibinfo{title}{The non-power model of the genetic code: a paradigm
  for interpreting genomic information}.
\newblock \bibinfo{journal}{Philosophical Transactions of the Royal Society A:
  Mathematical, Physical and Engineering Sciences} \bibinfo{volume}{374},
  \bibinfo{pages}{20150062}.
\newblock \DOIprefix\doi{10.1098/rsta.2015.0062}. \bibinfo{note}{non-power
  binary encoding of the genetic code using redundant signed-binary numeration
  with Fibonacci-like positional weights [1,1,2,4,7,8]; a strictly larger
  encoding family than the 24 power (Cartesian-product) bijections swept in
  this work.}
%Type = Article
\bibitem[{Grantham(1974)}]{grantham1974}
\bibinfo{author}{Grantham, R.}, \bibinfo{year}{1974}.
\newblock \bibinfo{title}{Amino acid difference formula to help explain protein
  evolution}.
\newblock \bibinfo{journal}{Science} \bibinfo{volume}{185},
  \bibinfo{pages}{862--864}.
\newblock \DOIprefix\doi{10.1126/science.185.4154.862}.
%Type = Article
\bibitem[{Grome et~al.(2025)Grome, Nguyen, Moonan, Mohler, Gurara, Wang, Hemez,
  Stenton, Cao, Radford, Kornaj, Patel, Prome, Rogulina, Sozanski, Tordoff,
  Rinehart and Isaacs}]{grome2025ochre}
\bibinfo{author}{Grome, M.W.}, \bibinfo{author}{Nguyen, M.T.A.},
  \bibinfo{author}{Moonan, D.W.}, \bibinfo{author}{Mohler, K.},
  \bibinfo{author}{Gurara, K.}, \bibinfo{author}{Wang, S.},
  \bibinfo{author}{Hemez, C.}, \bibinfo{author}{Stenton, B.J.},
  \bibinfo{author}{Cao, Y.}, \bibinfo{author}{Radford, F.},
  \bibinfo{author}{Kornaj, M.}, \bibinfo{author}{Patel, J.},
  \bibinfo{author}{Prome, M.}, \bibinfo{author}{Rogulina, S.},
  \bibinfo{author}{Sozanski, D.}, \bibinfo{author}{Tordoff, J.},
  \bibinfo{author}{Rinehart, J.}, \bibinfo{author}{Isaacs, F.J.},
  \bibinfo{year}{2025}.
\newblock \bibinfo{title}{Engineering a genomically recoded organism with one
  stop codon}.
\newblock \bibinfo{journal}{Nature} \bibinfo{volume}{639},
  \bibinfo{pages}{512--521}.
\newblock \DOIprefix\doi{10.1038/s41586-024-08501-x}.
%Type = Article
\bibitem[{Haig and Hurst(1991)}]{haig1991}
\bibinfo{author}{Haig, D.}, \bibinfo{author}{Hurst, L.D.},
  \bibinfo{year}{1991}.
\newblock \bibinfo{title}{A quantitative measure of error minimization in the
  genetic code}.
\newblock \bibinfo{journal}{Journal of Molecular Evolution}
  \bibinfo{volume}{33}, \bibinfo{pages}{412--417}.
\newblock \DOIprefix\doi{10.1007/BF02103132}.
%Type = Article
\bibitem[{Itzkovitz and Alon(2007)}]{itzkovitz2007}
\bibinfo{author}{Itzkovitz, S.}, \bibinfo{author}{Alon, U.},
  \bibinfo{year}{2007}.
\newblock \bibinfo{title}{The genetic code is nearly optimal for allowing
  additional information within protein-coding sequences}.
\newblock \bibinfo{journal}{Genome Research} \bibinfo{volume}{17},
  \bibinfo{pages}{405--412}.
\newblock \DOIprefix\doi{10.1101/gr.5987307}.
%Type = Article
\bibitem[{Jia and Jernigan(2021)}]{jia2021}
\bibinfo{author}{Jia, K.}, \bibinfo{author}{Jernigan, R.L.},
  \bibinfo{year}{2021}.
\newblock \bibinfo{title}{New amino acid substitution matrix brings sequence
  alignments into agreement with structure matches}.
\newblock \bibinfo{journal}{Proteins: Structure, Function, and Bioinformatics}
  \bibinfo{volume}{89}, \bibinfo{pages}{671--682}.
\newblock \DOIprefix\doi{10.1002/prot.26050}.
%Type = Article
\bibitem[{Kachale et~al.(2023)Kachale, Pavl{\'\i}kov{\'a}, Nenarokova,
  Roithov{\'a}, Durante, Milet{\'\i}nov{\'a}, Z{\'a}honov{\'a}, Nenarokov,
  Vot{\'y}pka, Hor{\'a}kov{\'a}, Ross, Yurchenko, Beznoskov{\'a}, Paris,
  Val{\'a}{\v s}ek and Luke{\v s}}]{kachale2023}
\bibinfo{author}{Kachale, A.}, \bibinfo{author}{Pavl{\'\i}kov{\'a}, Z.},
  \bibinfo{author}{Nenarokova, A.}, \bibinfo{author}{Roithov{\'a}, A.},
  \bibinfo{author}{Durante, I.M.}, \bibinfo{author}{Milet{\'\i}nov{\'a}, P.},
  \bibinfo{author}{Z{\'a}honov{\'a}, K.}, \bibinfo{author}{Nenarokov, S.},
  \bibinfo{author}{Vot{\'y}pka, J.}, \bibinfo{author}{Hor{\'a}kov{\'a}, E.},
  \bibinfo{author}{Ross, R.L.}, \bibinfo{author}{Yurchenko, V.},
  \bibinfo{author}{Beznoskov{\'a}, P.}, \bibinfo{author}{Paris, Z.},
  \bibinfo{author}{Val{\'a}{\v s}ek, L.S.}, \bibinfo{author}{Luke{\v s}, J.},
  \bibinfo{year}{2023}.
\newblock \bibinfo{title}{Short {tRNA} anticodon stem and mutant {eRF1} allow
  stop codon reassignment}.
\newblock \bibinfo{journal}{Nature} \bibinfo{volume}{613},
  \bibinfo{pages}{751--758}.
\newblock \DOIprefix\doi{10.1038/s41586-022-05584-2}.
%Type = Article
\bibitem[{Khrennikov and Kozyrev(2007)}]{khrennikov2007}
\bibinfo{author}{Khrennikov, A.Y.}, \bibinfo{author}{Kozyrev, S.V.},
  \bibinfo{year}{2007}.
\newblock \bibinfo{title}{Genetic code on the diadic plane}.
\newblock \bibinfo{journal}{Physica A: Statistical Mechanics and its
  Applications} \bibinfo{volume}{381}, \bibinfo{pages}{265--272}.
\newblock \DOIprefix\doi{10.1016/j.physa.2007.03.018}.
%Type = Article
\bibitem[{Koonin and Novozhilov(2009)}]{koonin2009}
\bibinfo{author}{Koonin, E.V.}, \bibinfo{author}{Novozhilov, A.S.},
  \bibinfo{year}{2009}.
\newblock \bibinfo{title}{Origin and evolution of the genetic code: the
  universal enigma}.
\newblock \bibinfo{journal}{IUBMB Life} \bibinfo{volume}{61},
  \bibinfo{pages}{99--111}.
\newblock \DOIprefix\doi{10.1002/iub.146}.
%Type = Article
\bibitem[{Kosiol et~al.(2004)Kosiol, Goldman and Buttimore}]{kosiol2004}
\bibinfo{author}{Kosiol, C.}, \bibinfo{author}{Goldman, N.},
  \bibinfo{author}{Buttimore, N.H.}, \bibinfo{year}{2004}.
\newblock \bibinfo{title}{A new criterion and method for amino acid
  classification}.
\newblock \bibinfo{journal}{Journal of Theoretical Biology}
  \bibinfo{volume}{228}, \bibinfo{pages}{97--106}.
\newblock \DOIprefix\doi{10.1016/j.jtbi.2003.12.010}.
%Type = Article
\bibitem[{Kyte and Doolittle(1982)}]{kyte1982}
\bibinfo{author}{Kyte, J.}, \bibinfo{author}{Doolittle, R.F.},
  \bibinfo{year}{1982}.
\newblock \bibinfo{title}{A simple method for displaying the hydropathic
  character of a protein}.
\newblock \bibinfo{journal}{Journal of Molecular Biology}
  \bibinfo{volume}{157}, \bibinfo{pages}{105--132}.
\newblock \DOIprefix\doi{10.1016/0022-2836(82)90515-0}.
%Type = Article
\bibitem[{Lajoie et~al.(2013)Lajoie, Rovner, Goodman, Aerni, Haimovich,
  Kuznetsov, Mercer, Wang, Carr, Mosberg, Rohland, Schultz, Jacobson, Rinehart,
  Church and Isaacs}]{lajoie2013}
\bibinfo{author}{Lajoie, M.J.}, \bibinfo{author}{Rovner, A.J.},
  \bibinfo{author}{Goodman, D.B.}, \bibinfo{author}{Aerni, H.R.},
  \bibinfo{author}{Haimovich, A.D.}, \bibinfo{author}{Kuznetsov, G.},
  \bibinfo{author}{Mercer, J.A.}, \bibinfo{author}{Wang, H.H.},
  \bibinfo{author}{Carr, P.A.}, \bibinfo{author}{Mosberg, J.A.},
  \bibinfo{author}{Rohland, N.}, \bibinfo{author}{Schultz, P.G.},
  \bibinfo{author}{Jacobson, J.M.}, \bibinfo{author}{Rinehart, J.},
  \bibinfo{author}{Church, G.M.}, \bibinfo{author}{Isaacs, F.J.},
  \bibinfo{year}{2013}.
\newblock \bibinfo{title}{Genomically recoded organisms expand biological
  functions}.
\newblock \bibinfo{journal}{Science} \bibinfo{volume}{342},
  \bibinfo{pages}{357--360}.
\newblock \DOIprefix\doi{10.1126/science.1241459}.
%Type = Article
\bibitem[{Miyata et~al.(1979)Miyata, Miyazawa and Yasunaga}]{miyata1979}
\bibinfo{author}{Miyata, T.}, \bibinfo{author}{Miyazawa, S.},
  \bibinfo{author}{Yasunaga, T.}, \bibinfo{year}{1979}.
\newblock \bibinfo{title}{Two types of amino acid substitutions in protein
  evolution}.
\newblock \bibinfo{journal}{Journal of Molecular Evolution}
  \bibinfo{volume}{12}, \bibinfo{pages}{219--236}.
\newblock \DOIprefix\doi{10.1007/BF01732340}.
%Type = Article
\bibitem[{Napolitano et~al.(2016)Napolitano, Landon, Gregg, Lajoie,
  Govindarajan, Mosberg, Kuznetsov, Goodman, Vargas-Rodriguez, Isaacs, S{\"o}ll
  and Church}]{napolitano2016}
\bibinfo{author}{Napolitano, M.G.}, \bibinfo{author}{Landon, M.},
  \bibinfo{author}{Gregg, C.J.}, \bibinfo{author}{Lajoie, M.J.},
  \bibinfo{author}{Govindarajan, L.}, \bibinfo{author}{Mosberg, J.A.},
  \bibinfo{author}{Kuznetsov, G.}, \bibinfo{author}{Goodman, D.B.},
  \bibinfo{author}{Vargas-Rodriguez, O.}, \bibinfo{author}{Isaacs, F.J.},
  \bibinfo{author}{S{\"o}ll, D.}, \bibinfo{author}{Church, G.M.},
  \bibinfo{year}{2016}.
\newblock \bibinfo{title}{Emergent rules for codon choice elucidated by editing
  rare arginine codons in {Escherichia coli}}.
\newblock \bibinfo{journal}{Proceedings of the National Academy of Sciences}
  \bibinfo{volume}{113}, \bibinfo{pages}{E5588--E5597}.
\newblock \DOIprefix\doi{10.1073/pnas.1605856113}.
%Type = Article
\bibitem[{Nyerges et~al.(2024)Nyerges, Chiappino-Pepe, Budnik, Baas-Thomas,
  Flynn, Yan, Ostrov, Liu, Wang, Zheng, Hu, Chen, Rudolph, Chen, Ahn, Spencer,
  Ayalavarapu, Tarver, Harmon-Smith, Hamilton, Blaby, Yoshikuni, Hajian, Jin,
  Kintses, Szamel, Seregi, Shen, Li and Church}]{nyerges2024}
\bibinfo{author}{Nyerges, {\'A}.}, \bibinfo{author}{Chiappino-Pepe, A.},
  \bibinfo{author}{Budnik, B.}, \bibinfo{author}{Baas-Thomas, M.},
  \bibinfo{author}{Flynn, R.}, \bibinfo{author}{Yan, S.},
  \bibinfo{author}{Ostrov, N.}, \bibinfo{author}{Liu, M.},
  \bibinfo{author}{Wang, M.}, \bibinfo{author}{Zheng, Q.}, \bibinfo{author}{Hu,
  F.}, \bibinfo{author}{Chen, K.}, \bibinfo{author}{Rudolph, A.},
  \bibinfo{author}{Chen, D.}, \bibinfo{author}{Ahn, J.},
  \bibinfo{author}{Spencer, O.}, \bibinfo{author}{Ayalavarapu, V.},
  \bibinfo{author}{Tarver, A.}, \bibinfo{author}{Harmon-Smith, M.},
  \bibinfo{author}{Hamilton, M.}, \bibinfo{author}{Blaby, I.},
  \bibinfo{author}{Yoshikuni, Y.}, \bibinfo{author}{Hajian, B.},
  \bibinfo{author}{Jin, A.}, \bibinfo{author}{Kintses, B.},
  \bibinfo{author}{Szamel, M.}, \bibinfo{author}{Seregi, V.},
  \bibinfo{author}{Shen, Y.}, \bibinfo{author}{Li, Z.},
  \bibinfo{author}{Church, G.M.}, \bibinfo{year}{2024}.
\newblock \bibinfo{title}{Synthetic genomes unveil the effects of synonymous
  recoding}.
\newblock \bibinfo{journal}{bioRxiv} ,
  \bibinfo{pages}{2024.06.16.599206}\DOIprefix\doi{10.1101/2024.06.16.599206}.
  \bibinfo{note}{preprint}.
%Type = Article
\bibitem[{Ostrov et~al.(2016)Ostrov, Landon, Guell, Kuznetsov, Teramoto,
  Cervantes, Zhou, Singh, Napolitano, Moosburner, Shrock, Pruitt, Conway,
  Goodman, Gardner, Tyree, Gonzales, Wanner, Norville, Lajoie and
  Church}]{ostrov2016}
\bibinfo{author}{Ostrov, N.}, \bibinfo{author}{Landon, M.},
  \bibinfo{author}{Guell, M.}, \bibinfo{author}{Kuznetsov, G.},
  \bibinfo{author}{Teramoto, J.}, \bibinfo{author}{Cervantes, N.},
  \bibinfo{author}{Zhou, M.}, \bibinfo{author}{Singh, K.},
  \bibinfo{author}{Napolitano, M.G.}, \bibinfo{author}{Moosburner, M.},
  \bibinfo{author}{Shrock, E.}, \bibinfo{author}{Pruitt, B.W.},
  \bibinfo{author}{Conway, N.}, \bibinfo{author}{Goodman, D.B.},
  \bibinfo{author}{Gardner, C.L.}, \bibinfo{author}{Tyree, G.},
  \bibinfo{author}{Gonzales, A.}, \bibinfo{author}{Wanner, B.L.},
  \bibinfo{author}{Norville, J.E.}, \bibinfo{author}{Lajoie, M.J.},
  \bibinfo{author}{Church, G.M.}, \bibinfo{year}{2016}.
\newblock \bibinfo{title}{Design, synthesis, and testing toward a 57-codon
  genome}.
\newblock \bibinfo{journal}{Science} \bibinfo{volume}{353},
  \bibinfo{pages}{819--822}.
\newblock \DOIprefix\doi{10.1126/science.aaf3639}.
%Type = Article
\bibitem[{Robertson et~al.(2025)Robertson, Rehm, Spinck, Schumann, Tian, Liu,
  Gu, Kleefeldt, Day, Liu, Christova, Z{\"u}rcher, B{\"o}ge, Birnbaum, van
  Bijsterveldt and Chin}]{robertson2025syn57}
\bibinfo{author}{Robertson, W.E.}, \bibinfo{author}{Rehm, F.B.H.},
  \bibinfo{author}{Spinck, M.}, \bibinfo{author}{Schumann, R.L.},
  \bibinfo{author}{Tian, R.}, \bibinfo{author}{Liu, W.}, \bibinfo{author}{Gu,
  Y.}, \bibinfo{author}{Kleefeldt, A.A.}, \bibinfo{author}{Day, C.F.},
  \bibinfo{author}{Liu, K.C.}, \bibinfo{author}{Christova, Y.},
  \bibinfo{author}{Z{\"u}rcher, J.F.}, \bibinfo{author}{B{\"o}ge, F.L.},
  \bibinfo{author}{Birnbaum, J.}, \bibinfo{author}{van Bijsterveldt, L.},
  \bibinfo{author}{Chin, J.W.}, \bibinfo{year}{2025}.
\newblock \bibinfo{title}{{Escherichia coli} with a 57-codon genetic code}.
\newblock \bibinfo{journal}{Science} \bibinfo{volume}{390},
  \bibinfo{pages}{eady4368}.
\newblock \DOIprefix\doi{10.1126/science.ady4368}.
%Type = Article
\bibitem[{Sengupta et~al.(2007)Sengupta, Yang and Higgs}]{sengupta2007}
\bibinfo{author}{Sengupta, S.}, \bibinfo{author}{Yang, X.},
  \bibinfo{author}{Higgs, P.G.}, \bibinfo{year}{2007}.
\newblock \bibinfo{title}{The mechanisms of codon reassignments in
  mitochondrial genetic codes}.
\newblock \bibinfo{journal}{Journal of Molecular Evolution}
  \bibinfo{volume}{64}, \bibinfo{pages}{662--688}.
\newblock \DOIprefix\doi{10.1007/s00239-006-0284-7}.
%Type = Article
\bibitem[{Singh et~al.(2023)Singh, Seah, Emmerich, Singh, Woehle, Huettel,
  Byerly, Stover, Sugiura, Harumoto and Swart}]{singh2023}
\bibinfo{author}{Singh, M.}, \bibinfo{author}{Seah, B.K.B.},
  \bibinfo{author}{Emmerich, C.}, \bibinfo{author}{Singh, A.},
  \bibinfo{author}{Woehle, C.}, \bibinfo{author}{Huettel, B.},
  \bibinfo{author}{Byerly, A.}, \bibinfo{author}{Stover, N.A.},
  \bibinfo{author}{Sugiura, M.}, \bibinfo{author}{Harumoto, T.},
  \bibinfo{author}{Swart, E.C.}, \bibinfo{year}{2023}.
\newblock \bibinfo{title}{Origins of genome-editing excisases as illuminated by
  the somatic genome of the ciliate {Blepharisma}}.
\newblock \bibinfo{journal}{Proceedings of the National Academy of Sciences}
  \bibinfo{volume}{120}, \bibinfo{pages}{e2213887120}.
\newblock \DOIprefix\doi{10.1073/pnas.2213887120}. \bibinfo{note}{reports the
  \emph{B.~stoltei} ATCC 30299 MAC genome assembly (GCA\_965603825.1; ENA
  Bstoltei\_ATCC30299\_MAC\_1.00) and documents that this strain uses an
  alternative nuclear genetic code with UGA codons reassigned from stops to
  tryptophan (main text and SI Appendix Fig.~S2B), in addition to the canonical
  NCBI table-15 UAG $\rightarrow$ Gln assignment.}
%Type = Article
\bibitem[{Su et~al.(2011)Su, Lieberman, Lang, Simonovi{\'c}, S{\"o}ll and
  Ling}]{su2011}
\bibinfo{author}{Su, D.}, \bibinfo{author}{Lieberman, A.},
  \bibinfo{author}{Lang, B.F.}, \bibinfo{author}{Simonovi{\'c}, M.},
  \bibinfo{author}{S{\"o}ll, D.}, \bibinfo{author}{Ling, J.},
  \bibinfo{year}{2011}.
\newblock \bibinfo{title}{An unusual trna-thr derived from trna-his reassigns
  in yeast mitochondria the {CUN} codons to threonine}.
\newblock \bibinfo{journal}{Nucleic Acids Research} \bibinfo{volume}{39},
  \bibinfo{pages}{4866--4874}.
\newblock \DOIprefix\doi{10.1093/nar/gkr073}.
%Type = Article
\bibitem[{Tsour et~al.(2026)Tsour, Machn{\'e}, Leduc, Widmer, Koo, Guez,
  Karczewski and Slavov}]{tsour2026}
\bibinfo{author}{Tsour, S.}, \bibinfo{author}{Machn{\'e}, R.},
  \bibinfo{author}{Leduc, A.}, \bibinfo{author}{Widmer, S.},
  \bibinfo{author}{Koo, E.}, \bibinfo{author}{Guez, J.},
  \bibinfo{author}{Karczewski, K.J.}, \bibinfo{author}{Slavov, N.},
  \bibinfo{year}{2026}.
\newblock \bibinfo{title}{Alternate {RNA} decoding results in stable and
  abundant proteins in mammals}.
\newblock \bibinfo{journal}{Nature} \DOIprefix\doi{10.1038/s41586-026-10678-2}.
%Type = Article
\bibitem[{Woese(1965)}]{woese1965}
\bibinfo{author}{Woese, C.R.}, \bibinfo{year}{1965}.
\newblock \bibinfo{title}{On the evolution of the genetic code}.
\newblock \bibinfo{journal}{Proceedings of the National Academy of Sciences}
  \bibinfo{volume}{54}, \bibinfo{pages}{1546--1552}.
\newblock \DOIprefix\doi{10.1073/pnas.54.6.1546}.
%Type = Article
\bibitem[{Woese(1973)}]{woese1973}
\bibinfo{author}{Woese, C.R.}, \bibinfo{year}{1973}.
\newblock \bibinfo{title}{Evolution of the genetic code}.
\newblock \bibinfo{journal}{Naturwissenschaften} \bibinfo{volume}{60},
  \bibinfo{pages}{447--459}.
\newblock \DOIprefix\doi{10.1007/BF00592854}.
%Type = Article
\bibitem[{Woese et~al.(1966)Woese, Dugre, Dugre, Kondo and
  Saxinger}]{woese1966}
\bibinfo{author}{Woese, C.R.}, \bibinfo{author}{Dugre, D.H.},
  \bibinfo{author}{Dugre, S.A.}, \bibinfo{author}{Kondo, M.},
  \bibinfo{author}{Saxinger, W.C.}, \bibinfo{year}{1966}.
\newblock \bibinfo{title}{On the fundamental nature and evolution of the
  genetic code}.
\newblock \bibinfo{journal}{Cold Spring Harbor Symposia on Quantitative
  Biology} \bibinfo{volume}{31}, \bibinfo{pages}{723--736}.
\newblock \DOIprefix\doi{10.1101/SQB.1966.031.01.093}.
%Type = Article
\bibitem[{Yurova~Axelsson and Khrennikov(2024a)}]{axelsson2024a}
\bibinfo{author}{Yurova~Axelsson, E.}, \bibinfo{author}{Khrennikov, A.Y.},
  \bibinfo{year}{2024}a.
\newblock \bibinfo{title}{Generation of genetic codes with 2-adic codon algebra
  and adaptive dynamics}.
\newblock \bibinfo{journal}{BioSystems} \bibinfo{volume}{240},
  \bibinfo{pages}{105230}.
\newblock \DOIprefix\doi{10.1016/j.biosystems.2024.105230}.
%Type = Article
\bibitem[{Yurova~Axelsson and Khrennikov(2024b)}]{axelsson2024b}
\bibinfo{author}{Yurova~Axelsson, E.}, \bibinfo{author}{Khrennikov, A.Y.},
  \bibinfo{year}{2024}b.
\newblock \bibinfo{title}{Universal dynamical function behind all genetic
  codes: {P}-adic attractor dynamical model}.
\newblock \bibinfo{journal}{BioSystems} \bibinfo{volume}{246},
  \bibinfo{pages}{105353}.
\newblock \DOIprefix\doi{10.1016/j.biosystems.2024.105353}.
%Type = Article
\bibitem[{Zehir et~al.(2017)Zehir, Benayed, Shah, Syed, Middha, Kim,
  Srinivasan, Gao, Chakravarty, Devlin, Hellmann, Barron, Schram, Hameed,
  Dogan, Ross, Hechtman, DeLair, Yao, Mandelker, Cheng, Chandramohan, Mohanty,
  Ptashkin, Jayakumaran, Prasad, Syed, Rema, Liu, Nafa, Borsu, Sadowska,
  Casanova, Bacares, Kiecka, Razumova, Son, Stewart, Baldi, Mullaney,
  Al-Ahmadie, Vakiani, Abeshouse, Penson, Jonsson, Camacho, Chang, Won, Gross,
  Kundra, Heins, Chen, Phillips, Zhang, Wang, Ochoa, Wills, Eubank, Thomas,
  Gardos, Reales, Galle, Durany, Cambria, Abida, Cercek, Feldman, Gounder,
  Hakimi, Harding, Iyer, Janjigian, Jordan, Kelly, Lowery, Morris, Omuro, Raj,
  Razavi, Shoushtari, Shukla, Soumerai, Varghese, Yaeger, Coleman, Bochner,
  Riely, Saltz, Scher, Sabbatini, Robson, Klimstra, Taylor, Baselga, Schultz,
  Hyman, Arcila, Solit, Ladanyi and Berger}]{zehir2017}
\bibinfo{author}{Zehir, A.}, \bibinfo{author}{Benayed, R.},
  \bibinfo{author}{Shah, R.H.}, \bibinfo{author}{Syed, A.},
  \bibinfo{author}{Middha, S.}, \bibinfo{author}{Kim, H.R.},
  \bibinfo{author}{Srinivasan, P.}, \bibinfo{author}{Gao, J.},
  \bibinfo{author}{Chakravarty, D.}, \bibinfo{author}{Devlin, S.M.},
  \bibinfo{author}{Hellmann, M.D.}, \bibinfo{author}{Barron, D.A.},
  \bibinfo{author}{Schram, A.M.}, \bibinfo{author}{Hameed, M.},
  \bibinfo{author}{Dogan, S.}, \bibinfo{author}{Ross, D.S.},
  \bibinfo{author}{Hechtman, J.F.}, \bibinfo{author}{DeLair, D.F.},
  \bibinfo{author}{Yao, J.}, \bibinfo{author}{Mandelker, D.L.},
  \bibinfo{author}{Cheng, D.T.}, \bibinfo{author}{Chandramohan, R.},
  \bibinfo{author}{Mohanty, A.S.}, \bibinfo{author}{Ptashkin, R.N.},
  \bibinfo{author}{Jayakumaran, G.}, \bibinfo{author}{Prasad, M.},
  \bibinfo{author}{Syed, M.H.}, \bibinfo{author}{Rema, A.B.},
  \bibinfo{author}{Liu, Z.Y.}, \bibinfo{author}{Nafa, K.},
  \bibinfo{author}{Borsu, L.}, \bibinfo{author}{Sadowska, J.},
  \bibinfo{author}{Casanova, J.}, \bibinfo{author}{Bacares, R.},
  \bibinfo{author}{Kiecka, I.J.}, \bibinfo{author}{Razumova, A.},
  \bibinfo{author}{Son, J.B.}, \bibinfo{author}{Stewart, L.},
  \bibinfo{author}{Baldi, T.}, \bibinfo{author}{Mullaney, K.A.},
  \bibinfo{author}{Al-Ahmadie, H.}, \bibinfo{author}{Vakiani, E.},
  \bibinfo{author}{Abeshouse, A.A.}, \bibinfo{author}{Penson, A.V.},
  \bibinfo{author}{Jonsson, P.}, \bibinfo{author}{Camacho, N.},
  \bibinfo{author}{Chang, M.T.}, \bibinfo{author}{Won, H.H.},
  \bibinfo{author}{Gross, B.E.}, \bibinfo{author}{Kundra, R.},
  \bibinfo{author}{Heins, Z.J.}, \bibinfo{author}{Chen, H.W.},
  \bibinfo{author}{Phillips, S.}, \bibinfo{author}{Zhang, H.},
  \bibinfo{author}{Wang, J.}, \bibinfo{author}{Ochoa, A.},
  \bibinfo{author}{Wills, J.}, \bibinfo{author}{Eubank, M.},
  \bibinfo{author}{Thomas, S.B.}, \bibinfo{author}{Gardos, S.M.},
  \bibinfo{author}{Reales, D.N.}, \bibinfo{author}{Galle, J.},
  \bibinfo{author}{Durany, R.}, \bibinfo{author}{Cambria, R.},
  \bibinfo{author}{Abida, W.}, \bibinfo{author}{Cercek, A.},
  \bibinfo{author}{Feldman, D.R.}, \bibinfo{author}{Gounder, M.M.},
  \bibinfo{author}{Hakimi, A.A.}, \bibinfo{author}{Harding, J.J.},
  \bibinfo{author}{Iyer, G.}, \bibinfo{author}{Janjigian, Y.Y.},
  \bibinfo{author}{Jordan, E.J.}, \bibinfo{author}{Kelly, C.M.},
  \bibinfo{author}{Lowery, M.A.}, \bibinfo{author}{Morris, L.G.T.},
  \bibinfo{author}{Omuro, A.M.}, \bibinfo{author}{Raj, N.},
  \bibinfo{author}{Razavi, P.}, \bibinfo{author}{Shoushtari, A.N.},
  \bibinfo{author}{Shukla, N.}, \bibinfo{author}{Soumerai, T.E.},
  \bibinfo{author}{Varghese, A.M.}, \bibinfo{author}{Yaeger, R.},
  \bibinfo{author}{Coleman, J.}, \bibinfo{author}{Bochner, B.},
  \bibinfo{author}{Riely, G.J.}, \bibinfo{author}{Saltz, L.B.},
  \bibinfo{author}{Scher, H.I.}, \bibinfo{author}{Sabbatini, P.J.},
  \bibinfo{author}{Robson, M.E.}, \bibinfo{author}{Klimstra, D.S.},
  \bibinfo{author}{Taylor, B.S.}, \bibinfo{author}{Baselga, J.},
  \bibinfo{author}{Schultz, N.}, \bibinfo{author}{Hyman, D.M.},
  \bibinfo{author}{Arcila, M.E.}, \bibinfo{author}{Solit, D.B.},
  \bibinfo{author}{Ladanyi, M.}, \bibinfo{author}{Berger, M.F.},
  \bibinfo{year}{2017}.
\newblock \bibinfo{title}{Mutational landscape of metastatic cancer revealed
  from prospective clinical sequencing of 10,000 patients}.
\newblock \bibinfo{journal}{Nature Medicine} \bibinfo{volume}{23},
  \bibinfo{pages}{703--713}.
\newblock \DOIprefix\doi{10.1038/nm.4333}.

\end{thebibliography}


\end{document}
